\section{Haag Kastler Axioms}\label{sctn:haag-kastler-axioms}
With that out of the way we can state the axioms. Each axiom below is presented as a \emph{definition} so that it is captured as a node in the blueprint declaration graph. In the Lean formalization each axiom corresponds to a \texttt{Prop}-valued predicate (or, in the case of \emph{Local Algebras}, a data field), and the bundling structure \texttt{HaagKastlerNet} packages them together (see \ref{def:haag-kastler-net}). Downstream theorems take an instance of \texttt{HaagKastlerNet} as a hypothesis and invoke each axiom as a projection.

\begin{definition}[Axiom 1: Local Algebras]
    \label{def:local-algebras}
    \lean{Physicslib4.AQFT.HaagKastler.LocalNet}
    \leanfile{Physicslib4/AQFT/HaagKastler/LocalAlgebras.lean}
    \leanok
    \uses{def:alexandrov-topology, def:minkowski-spacetime, def:chronological-future-and-chronological-past}
    For any basis element $\mathbf{B}$ of the Alexandrov topology on Minkowski spacetime, i.e. any set of the form $I^+(p) \cap I^-(q)$, there is a corresponding abstract C*-algebra $\mathfrak{U}(\mathbf{B})$
    \begin{align}
        \mathbf{B} \longmapsto \mathfrak{U}(\mathbf{B}),
    \end{align}
    and when $\mathbf{B}$ is the empty set, we have the distinguished correspondence
    \begin{align}
        \emptyset \longmapsto \mathbb{C} \mathbf{1},
    \end{align}
    where $\mathbf{1}$ is the multiplicative identity in the abstract C*-algebra $\mathbb{C} \mathbf{1}$.
\end{definition}

\begin{definition}[Axiom 2: Isotony]
    \label{def:isotony}
    \lean{Physicslib4.AQFT.HaagKastler.Isotony}
    \leanfile{Physicslib4/AQFT/HaagKastler/Isotony.lean}
    \leanok
    \uses{def:alexandrov-topology, def:minkowski-spacetime, def:chronological-future-and-chronological-past, def:local-algebras}
    Let $\mathbf{B}_1$ and $\mathbf{B}_2$ be any two basis elements of the Alexandrov topology on Minkowski spacetime, i.e. any two sets of the form $I^+(p_1) \cap I^-(q_1)$ and $I^+(p_2) \cap I^-(q_2)$.

    The axiom supplies, as \emph{data}, a family of unital $*$-monomorphisms
    \begin{align}
        i_{\mathbf{B}_1\mathbf{B}_2} : \mathfrak{U}(\mathbf{B}_1) \hookrightarrow \mathfrak{U}(\mathbf{B}_2),
    \end{align}
    one for every pair of basis sets with $\mathbf{B}_1 \subseteq \mathbf{B}_2$, subject to three conditions:
    \begin{itemize}
        \item[(a)] \emph{injectivity}: each $i_{\mathbf{B}_1\mathbf{B}_2}$ is injective;
        \item[(b)] \emph{identity}: $i_{\mathbf{B}\mathbf{B}} = \mathrm{id}_{\mathfrak{U}(\mathbf{B})}$ for every basis set $\mathbf{B}$;
        \item[(c)] \emph{composition}: $i_{\mathbf{B}_2\mathbf{B}_3} \circ i_{\mathbf{B}_1\mathbf{B}_2} = i_{\mathbf{B}_1\mathbf{B}_3}$ whenever $\mathbf{B}_1 \subseteq \mathbf{B}_2 \subseteq \mathbf{B}_3$.
    \end{itemize}
    Conditions (b) and (c) say exactly that $\mathbf{B} \mapsto \mathfrak{U}(\mathbf{B})$ is a \emph{functor} on the inclusion order of basis sets, with the $i_{\mathbf{B}_1\mathbf{B}_2}$ as its action on morphisms.

    Two points about the shape of this axiom. First, the family must be \emph{chosen data} and not an existence statement: (b) and (c) are equations between the maps themselves, so there is nothing to state unless the maps are fixed. An axiom of the form ``for each inclusion there exists some monomorphism'' cannot express functoriality at all.

    Second, the hypothesis is the \emph{non-strict} inclusion $\mathbf{B}_1 \subseteq \mathbf{B}_2$, matching the formalisation, which quantifies over \texttt{B\ensuremath{_1} \ensuremath{\subseteq} B\ensuremath{_2}}. Earlier versions of this statement wrote $\subset$; that was a divergence from the Lean, and the non-strict form is the correct one. It also matters mathematically: the reflexive case $\mathbf{B}_1 = \mathbf{B}_2$ is what makes (b) expressible, and the strict reading would leave the diagonal of the inclusion order outside the axiom altogether.

    Conditions (b) and (c) are required rather than derived because they are properties of the net's \emph{chosen} embeddings, not of the spacetime: no geometric fact about Alexandrov diamonds constrains which monomorphism a net picks for a given inclusion, so coherence cannot be discharged after the fact and must be part of the axiom. Their payoff is that the assignment becomes a genuine directed system, which is what gives \ref{def:quasilocal-algebra} its algebra structure, and that consumers no longer have to carry coherence as a side hypothesis.
\end{definition}

\subsection{The Quasilocal Colimit}

The quasilocal algebra of \ref{def:quasilocal-algebra} below rests on a chain of algebraic facts about the directed system of Axiom 2, which we record in this subsection. They fall into two groups: those putting a normed $*$-algebra structure on the colimit, assembled in \ref{lmm:quasilocal-union-normed-star-algebra}, and those passing from that to its completion, assembled in \ref{lmm:quasilocal-completion-cstar}. Axioms 3--5 depend only on \ref{def:quasilocal-algebra} itself and not on how it is built, so a reader after the axiom presentation alone may skip this subsection.

\begin{lemma}[Alexandrov Diamonds are Directed under Inclusion]
    \label{lmm:alexandrov-diamonds-isDirected}
    \lean{Physicslib4.AQFT.HaagKastler.Diamond, Physicslib4.AQFT.HaagKastler.directedOn_alexandrovBasis, Physicslib4.AQFT.HaagKastler.instIsDirectedOrderDiamond}
    \uses{def:alexandrov-topology, lmm:minkowski-diamonds-upward-directed}
    \leanok
    The order-theoretic repackaging of \ref{lmm:minkowski-diamonds-upward-directed}: the set of Alexandrov diamonds on standard Minkowski spacetime is \texttt{DirectedOn (\ensuremath{\subseteq})}, and hence, viewed as a subtype ordered by inclusion, satisfies \texttt{IsDirectedOrder}.
\end{lemma}
\begin{proof}
    \leanok
    \uses{lmm:minkowski-diamonds-upward-directed}
    Unfolding \texttt{DirectedOn} turns the goal into the statement of \ref{lmm:minkowski-diamonds-upward-directed}, which is therefore applied directly. The passage from the set-level \texttt{DirectedOn} to the subtype-level \texttt{IsDirectedOrder} that the colimit construction wants as an instance is Mathlib's \texttt{DirectedOn.isDirectedOrder}. There is no mathematical content here beyond \ref{lmm:minkowski-diamonds-upward-directed}; the node exists so that the colimit nodes can cite an instance rather than restate a $\forall\exists$ statement.
\end{proof}

\begin{lemma}[The Isotony Family is a Directed System]
    \label{lmm:isotony-directed-system}
    \lean{Physicslib4.AQFT.HaagKastler.transitionHom, Physicslib4.AQFT.HaagKastler.instDirectedSystemIsotony}
    \uses{def:local-algebras, def:isotony}
    \leanok
    The isotony family $i_{\mathbf{B}_1\mathbf{B}_2}$ of Axiom 2 (\ref{def:isotony}), indexed by the diamonds ordered by inclusion, is a \emph{directed system} in Mathlib's sense: it satisfies \texttt{DirectedSystem}.
\end{lemma}
\begin{proof}
    \leanok
    \uses{def:isotony}
    Mathlib's \texttt{DirectedSystem f} is a two-field class: \texttt{map\_self}, saying $f_{\le\mathrm{rfl}} = \mathrm{id}$, and \texttt{map\_map}, saying $f_{jk} \circ f_{ij} = f_{ik}$. These are literally Axiom 2(b) and Axiom 2(c) (\ref{def:isotony}), so the instance is constructed by supplying those two projections. This is the node that cashes in the remark in \ref{def:isotony} that (b) and (c) make the assignment a functor.
\end{proof}

\emph{The colimit and its algebraic structure come from Mathlib.} Let $\mathcal{D}$ be the type of Alexandrov diamonds ordered by inclusion and set
\begin{align}
    \varinjlim_{\mathbf{B}} \mathfrak{U}(\mathbf{B}) := \texttt{DirectLimit } \mathfrak{U}\ i,
\end{align}
Mathlib's general direct limit: the quotient of the sigma type $\Sigma_{\mathbf{B} : \mathcal{D}} \mathfrak{U}(\mathbf{B})$ by \texttt{DirectLimit.setoid}, which identifies $\langle \mathbf{B}_1, a_1\rangle$ with $\langle \mathbf{B}_2, a_2\rangle$ exactly when $i_{\mathbf{B}_1\mathbf{B}}(a_1) = i_{\mathbf{B}_2\mathbf{B}}(a_2)$ for some common $\mathbf{B}$. Two side conditions make this well posed and are both discharged above: the hypothesis that the $i$ form a \texttt{DirectedSystem} is \ref{lmm:isotony-directed-system}, and the hypothesis that the index order is directed --- needed already for \texttt{DirectLimit.setoid} to be transitive --- comes from \ref{lmm:alexandrov-diamonds-isDirected} through \texttt{DirectedOn.isDirectedOrder}.

There are no blueprint nodes for the algebraic structure on this object, because Mathlib already provides all of it, mostly by typeclass inference. Concretely:
\begin{itemize}
    \item \emph{Ring structure and canonical maps.} The \texttt{Ring (DirectLimit G f)} instance --- declared anonymously, so there is no declaration name to cite and it is found by typeclass inference --- together with $\iota_{\mathbf{B}} = \texttt{DirectLimit.Ring.of}$ and the universal property \texttt{DirectLimit.Ring.lift}. Concretely the product of $\iota_{\mathbf{B}_1}(a_1)$ and $\iota_{\mathbf{B}_2}(a_2)$ is $\iota_{\mathbf{B}}\bigl(i_{\mathbf{B}_1\mathbf{B}}(a_1)\, i_{\mathbf{B}_2\mathbf{B}}(a_2)\bigr)$ for any diamond $\mathbf{B}$ containing both, and functoriality is what makes this independent of the choice of $\mathbf{B}$.
    \item \emph{Involution and the $*$-ring axioms.} The \texttt{Star (DirectLimit G f)} instance, whose action on representatives is \texttt{DirectLimit.star\_def}: it sends the class of $\langle \mathbf{B}, a\rangle$ to the class of $\langle \mathbf{B}, a^{*}\rangle$, so $\iota_{\mathbf{B}}(a)^{*} = \iota_{\mathbf{B}}(a^{*})$ holds by definition. The identities $x^{**} = x$, $(x+y)^{*} = x^{*}+y^{*}$ and $(xy)^{*} = y^{*}x^{*}$ are the \texttt{StarRing (DirectLimit G f)} instance.
    \item \emph{$\mathbb{C}$-algebra structure.} The \texttt{Algebra R (DirectLimit G f)} instance --- again declared anonymously, with no declaration name to cite, and found by typeclass inference --- with canonical algebra maps \texttt{DirectLimit.Algebra.of}, so each $\iota_{\mathbf{B}}$ is a unital $\mathbb{C}$-algebra homomorphism. Compatibility of $\star$ with scalars, $(c \cdot x)^{*} = \bar{c} \cdot x^{*}$, is the \texttt{StarModule \ensuremath{\mathbb{C}} (DirectLimit G f)} instance (autogenerated name \texttt{DirectLimit.instStarModule}). Together with the previous item this is the entire $*$-algebra structure over $\mathbb{C}$; nothing has to be constructed by hand.
    \item \emph{Injectivity of the canonical maps.} Each $\iota_{\mathbf{B}}$ is injective by \texttt{DirectLimit.mk\_injective}, whose hypothesis --- injectivity of every transition map --- is exactly Axiom 2(a) (\ref{def:isotony}). This is where Axiom 2(a) earns its place: without it the colimit could identify distinct local observables, and both positive definiteness of the colimit norm below and the reading of $\mathfrak{U}(\mathbf{B})$ as a subalgebra of $\mathfrak{U}$ would fail.
\end{itemize}

\emph{Which Mathlib construction.} It must be the general \texttt{DirectLimit}, \emph{not} \texttt{Ring.DirectLimit}. Despite its name the latter is defined as a quotient of \texttt{FreeCommRing (\ensuremath{\Sigma} i, G i)} and its \texttt{deriving} clause produces a \texttt{CommRing}; it is a commutative-rings-only construction and cannot carry the local algebras, which are noncommutative in every case of interest. The general \texttt{DirectLimit} imposes no commutativity: it asks only that the transition maps be \texttt{RingHomClass}, which the unital $*$-monomorphisms of Axiom 2 are.

One more fact needs no node either: each isotony embedding $i_{\mathbf{B}_1\mathbf{B}_2}$ is \emph{isometric}. An injective $*$-homomorphism between complex C*-algebras is norm preserving, which is \texttt{NonUnitalStarAlgHom.norm\_map}, \texttt{lemma norm\_map (\ensuremath{\varphi} : F) (h\ensuremath{\varphi} : Function.Injective \ensuremath{\varphi}) (a : A) : \ensuremath{\|\varphi a\| = \|a\|}}. Its only hypothesis is injectivity, i.e.\ Axiom 2(a) (\ref{def:isotony}), and the local algebras are complex C*-algebras by Axiom 1 (\ref{def:local-algebras}); this repository already applies that citation twice.

What is left for this blueprint is therefore only the \emph{norm} on the colimit, which Mathlib cannot guess, and its passage to the completion.

\begin{lemma}[The Colimit Norm is Well Defined]
    \label{lmm:quasilocal-colimit-norm-well-defined}
    \lean{Physicslib4.AQFT.HaagKastler.QuasilocalColimit, Physicslib4.AQFT.HaagKastler.norm_transitionHom, Physicslib4.AQFT.HaagKastler.colimitNorm, Physicslib4.AQFT.HaagKastler.colimitNorm_mk}
    \uses{def:local-algebras, def:isotony, lmm:isotony-directed-system, lmm:alexandrov-diamonds-isDirected}
    \leanok
    Setting $\|[a]\| := \|a\|$ for a representative $a \in \mathfrak{U}(\mathbf{B})$ gives a well-defined function on the colimit $\varinjlim_{\mathbf{B}} \mathfrak{U}(\mathbf{B})$: the value depends neither on the diamond $\mathbf{B}$ nor on the representative $a$ chosen.
\end{lemma}
\begin{proof}
    \leanok
    \uses{def:local-algebras, def:isotony, lmm:alexandrov-diamonds-isDirected}
    Two representatives of the same class become equal after transport into a common containing diamond, which exists by \ref{lmm:alexandrov-diamonds-isDirected}. That transport is by isotony embeddings, and those are isometric --- \texttt{NonUnitalStarAlgHom.norm\_map} applied with Axiom 2(a) (\ref{def:isotony}) in the complex C*-algebras of Axiom 1 (\ref{def:local-algebras}) --- so each representative has the same norm as their common image and hence the same norm as the other.
\end{proof}

\begin{lemma}[Common Representatives for Two Colimit Elements]
    \label{lmm:quasilocal-colimit-common-representatives}
    \lean{Physicslib4.AQFT.HaagKastler.exists_common_representatives}
    \uses{def:local-algebras, def:isotony, lmm:isotony-directed-system, lmm:alexandrov-diamonds-isDirected}
    \leanok
    For any two elements $x, y$ of $\varinjlim_{\mathbf{B}} \mathfrak{U}(\mathbf{B})$ there are a single diamond $\mathbf{B}$ and elements $a, b \in \mathfrak{U}(\mathbf{B})$ with $x = \iota_{\mathbf{B}}(a)$ and $y = \iota_{\mathbf{B}}(b)$.
\end{lemma}
\begin{proof}
    \leanok
    \uses{lmm:alexandrov-diamonds-isDirected, lmm:isotony-directed-system}
    This is Mathlib's \texttt{DirectLimit.exists\_eq\_mk\ensuremath{_2}}, \texttt{theorem exists\_eq\_mk\ensuremath{_2} (z w : DirectLimit F f) : \ensuremath{\exists} i x y, z = [\ensuremath{\langle}i, x\ensuremath{\rangle}] \ensuremath{\wedge} w = [\ensuremath{\langle}i, y\ensuremath{\rangle}]}, applied to the directed system \ref{lmm:isotony-directed-system}; its directedness hypothesis is \ref{lmm:alexandrov-diamonds-isDirected} through \texttt{DirectedOn.isDirectedOrder}. When the goal is a proposition rather than data, the same fact is packaged as the induction principle \texttt{DirectLimit.induction\ensuremath{_2}}, which is the form actually used in tactic proofs.

    The node exists because every two-element claim about the colimit norm below opens with this move --- subadditivity and submultiplicativity in \ref{lmm:quasilocal-colimit-norm-axioms} --- and in Lean it is a two-line \texttt{obtain} that would otherwise be repeated verbatim.
\end{proof}

\begin{lemma}[The Colimit Norm is a Ring Norm and a Normed-Space Norm]
    \label{lmm:quasilocal-colimit-norm-axioms}
    \lean{Physicslib4.AQFT.HaagKastler.instNonemptyDiamond, Physicslib4.AQFT.HaagKastler.colimitRingNorm, Physicslib4.AQFT.HaagKastler.colimitNormedRing, Physicslib4.AQFT.HaagKastler.colimitNorm_smul}
    \uses{lmm:quasilocal-colimit-norm-well-defined, lmm:quasilocal-colimit-common-representatives, def:local-algebras, def:isotony, lmm:alexandrov-diamonds-isDirected}
    \leanok
    The norm of \ref{lmm:quasilocal-colimit-norm-well-defined} satisfies the normed-$*$-algebra norm axioms on $\varinjlim_{\mathbf{B}} \mathfrak{U}(\mathbf{B})$: for all $x, y$ in the colimit and all $c \in \mathbb{C}$,
    \begin{align}
        \|0\| = 0, \qquad \|{-x}\| = \|x\|, \qquad \|x + y\| \le \|x\| + \|y\|, \qquad \|xy\| \le \|x\|\,\|y\|,
    \end{align}
    \begin{align}
        \|x\| = 0 \implies x = 0, \qquad \|c \cdot x\| = |c|\,\|x\|.
    \end{align}
    The first five clauses are exactly the fields of a \texttt{RingNorm} on the colimit: \texttt{map\_zero'} and \texttt{neg'} are inherited from \texttt{AddGroupSeminorm} through \texttt{RingSeminorm}, \texttt{add\_le'} and \texttt{mul\_le'} are subadditivity and submultiplicativity, and \texttt{eq\_zero\_of\_map\_eq\_zero'} is positive definiteness. So proving them \emph{is} building the \texttt{RingNorm}, whence a \texttt{NormedRing} structure by \texttt{RingNorm.toNormedRing}. Absolute homogeneity is listed separately because it is \emph{not} a \texttt{RingNorm} field at all --- a \texttt{RingNorm} knows nothing about the scalars --- and its only role is to supply the \texttt{norm\_smul\_le} field of \texttt{NormedSpace \ensuremath{\mathbb{C}}} over the \texttt{NormedRing} structure just obtained.

    All six clauses are gathered into one node because in Lean they are the fields of a single \texttt{NormedRing} plus \texttt{NormedSpace \ensuremath{\mathbb{C}}} instance and each is about two lines of the same transport argument; separating them would be over-decomposition.

    The \texttt{RingNorm} detour is not bureaucracy: \texttt{NormedRing} bundles a \texttt{MetricSpace}, and there is no metric on the colimit quotient until this norm provides one, so \texttt{NormedRing} cannot even be stated first and then filled in field by field. The order is forced: build the bare function (\ref{lmm:quasilocal-colimit-norm-well-defined}), prove the \texttt{RingNorm} fields, and only then obtain \texttt{NormedRing} and \texttt{NormedSpace \ensuremath{\mathbb{C}}} from it.

    $*$-invariance $\|x^{*}\| = \|x\|$ is deliberately not listed. It is not an independent obligation: once the C*-inequality of \ref{lmm:quasilocal-colimit-cstar-identity} is available, \texttt{CStarRing.to\_normedStarGroup} produces the \texttt{NormedStarGroup} instance, and $\|x^{*}\| = \|x\|$ with it.
\end{lemma}
\begin{proof}
    \leanok
    \uses{lmm:quasilocal-colimit-norm-well-defined, lmm:quasilocal-colimit-common-representatives, def:local-algebras, def:isotony}
    Every clause is transport along a representative, which is why they share a node. For the two binary clauses take a common diamond $\mathbf{B}$ and representatives $a, b \in \mathfrak{U}(\mathbf{B})$ with $x = \iota_{\mathbf{B}}(a)$, $y = \iota_{\mathbf{B}}(b)$ by \ref{lmm:quasilocal-colimit-common-representatives}; for the unary clauses one representative suffices. Since $\iota_{\mathbf{B}}$ is a unital $\mathbb{C}$-algebra homomorphism (the \texttt{Algebra} instance recorded at the start of this subsection), $0$, $-x$, $c \cdot x$, $x + y$ and $xy$ are $\iota_{\mathbf{B}}(0)$, $\iota_{\mathbf{B}}(-a)$, $\iota_{\mathbf{B}}(c \cdot a)$, $\iota_{\mathbf{B}}(a+b)$ and $\iota_{\mathbf{B}}(ab)$.

    Rewriting every colimit norm as a $\mathfrak{U}(\mathbf{B})$-norm with \ref{lmm:quasilocal-colimit-norm-well-defined} then reduces the clauses to \texttt{norm\_zero}, \texttt{norm\_neg}, \texttt{norm\_add\_le}, \texttt{norm\_mul\_le} and \texttt{norm\_smul} in the C*-algebra $\mathfrak{U}(\mathbf{B})$ (\ref{def:local-algebras}).

    Positive definiteness is the one clause with content beyond transport, and it is where Axiom 2(a) is spent. From $\|x\| = 0$ and \ref{lmm:quasilocal-colimit-norm-well-defined} we get $\|a\| = 0$, hence $a = 0$ by \texttt{norm\_eq\_zero} in $\mathfrak{U}(\mathbf{B})$; the step from $a = 0$ to $x = 0$ is injectivity of $\iota_{\mathbf{B}}$, which is \texttt{DirectLimit.mk\_injective} applied with Axiom 2(a) (\ref{def:isotony}).

    Assembling the \texttt{RingNorm} and feeding it to \texttt{RingNorm.toNormedRing}, then adding homogeneity as \texttt{norm\_smul\_le}, is one term each.
\end{proof}

\begin{lemma}[The Quasilocal Union is a Normed $*$-Algebra]
    \label{lmm:quasilocal-union-normed-star-algebra}
    \lean{Physicslib4.AQFT.HaagKastler.norm_eq_colimitNorm, Physicslib4.AQFT.HaagKastler.colimitNormedAlgebra}
    \uses{def:local-algebras, def:isotony, lmm:minkowski-diamonds-upward-directed, lmm:isotony-directed-system, lmm:quasilocal-colimit-norm-well-defined, lmm:quasilocal-colimit-norm-axioms}
    \leanok
    The colimit of the local algebras $\mathfrak{U}(\mathbf{B})$ (\ref{def:local-algebras}) along the isotony family of Axiom 2 (\ref{def:isotony}), taken over the upward-directed Alexandrov diamonds (\ref{lmm:minkowski-diamonds-upward-directed}), carries a well-defined normed $*$-algebra structure over $\mathbb{C}$. Explicitly it supplies, on the colimit, a \texttt{NormedRing} structure, a \texttt{StarRing} structure, a \texttt{NormedAlgebra \ensuremath{\mathbb{C}}} structure and a \texttt{StarModule \ensuremath{\mathbb{C}}} structure.
\end{lemma}
\begin{proof}
    \leanok
    \uses{lmm:isotony-directed-system, lmm:quasilocal-colimit-norm-well-defined, lmm:quasilocal-colimit-norm-axioms}
    Assemble. The $*$-algebra structure over $\mathbb{C}$ is supplied by Mathlib's \texttt{DirectLimit} instances, recorded at the start of this subsection and found by typeclass inference once \ref{lmm:isotony-directed-system} is in place; \ref{lmm:quasilocal-colimit-norm-well-defined} supplies the norm function; \ref{lmm:quasilocal-colimit-norm-axioms} supplies the axioms that norm satisfies, in the assembled form \texttt{NormedRing} plus \texttt{NormedSpace \ensuremath{\mathbb{C}}}. The \texttt{NormedAlgebra \ensuremath{\mathbb{C}}} structure asserted above is exactly that \texttt{NormedSpace \ensuremath{\mathbb{C}}} together with the \texttt{Algebra \ensuremath{\mathbb{C}}} instance from \texttt{DirectLimit}: \texttt{NormedAlgebra} has no field beyond \texttt{Algebra} other than \texttt{norm\_smul\_le}, which is the \texttt{NormedSpace} field. It is stated in the \texttt{NormedAlgebra} form here because that is what \ref{def:completion-standing-hypotheses} consumes downstream. Bundling them is the normed $*$-algebra structure.

    The \texttt{StarModule \ensuremath{\mathbb{C}}} field, $(c \cdot x)^{*} = \bar{c} \cdot x^{*}$, is explicitly part of what this node supplies, and it costs nothing: it is the \texttt{StarModule \ensuremath{\mathbb{C}} (DirectLimit G f)} instance (autogenerated name \texttt{DirectLimit.instStarModule}), found by typeclass inference from the corresponding \texttt{StarModule \ensuremath{\mathbb{C}}} structure on each local algebra. It is listed because downstream it is a required field of \texttt{CStarAlgebra}, and on the completion side (\ref{lmm:star-extends-to-completion}) the analogous field is \emph{not} free.

    Nothing is claimed here about the involution being isometric, and nothing needs to be: $\|x^{*}\| = \|x\|$ comes free from the C*-inequality (\ref{lmm:quasilocal-colimit-cstar-identity}) through \texttt{CStarRing.to\_normedStarGroup}.
\end{proof}

\begin{lemma}[The Colimit Satisfies the C*-Inequality]
    \label{lmm:quasilocal-colimit-cstar-identity}
    \lean{Physicslib4.AQFT.HaagKastler.colimitCStarRing}
    \uses{lmm:quasilocal-union-normed-star-algebra, lmm:quasilocal-colimit-norm-well-defined, def:local-algebras}
    \leanok
    The normed $*$-algebra of \ref{lmm:quasilocal-union-normed-star-algebra} satisfies
    \begin{align}
        \|x\| \, \|x\| \le \|x^{*} x\|
    \end{align}
    for every $x$ in the colimit.

    Only this inequality is asserted, because it is literally the single field \texttt{norm\_mul\_self\_le} of Mathlib's \texttt{CStarRing} class, so establishing it \emph{is} establishing the \texttt{CStarRing} instance. The familiar equality $\|x^{*}x\| = \|x\|^{2}$ then comes back for free from \texttt{CStarRing.norm\_star\_mul\_self}, and $*$-invariance $\|x^{*}\| = \|x\|$ from \texttt{CStarRing.to\_normedStarGroup}; neither needs a proof of its own. The colimit is still not a C*-algebra: it is in general not complete.
\end{lemma}
\begin{proof}
    \leanok
    \uses{lmm:quasilocal-colimit-norm-well-defined, def:local-algebras}
    A single element is involved, so no directedness is needed here. Pick a representative $a \in \mathfrak{U}(\mathbf{B})$ of $x$. Then $x^{*}x$ has the representative $a^{*}a$ in the same $\mathfrak{U}(\mathbf{B})$, since the involution acts on representatives by \texttt{DirectLimit.star\_def}, and by \ref{lmm:quasilocal-colimit-norm-well-defined} the colimit norms of $x$ and $x^{*}x$ are the norms of $a$ and $a^{*}a$ there. So the claim reduces to $\|a\|\,\|a\| \le \|a^{*}a\|$ in $\mathfrak{U}(\mathbf{B})$, which is the \texttt{CStarRing} instance carried by the local algebras (\ref{def:local-algebras}), available as \texttt{CStarRing.norm\_mul\_self\_le}.
\end{proof}

The remaining results of this subsection have nothing to do with the quasilocal setting and are stated for an arbitrary normed $*$-algebra, as their titles already claim; they are applied to the colimit only in \ref{lmm:quasilocal-completion-cstar}. Their common hypotheses are collected once, as a node of their own so that each of them can cite it:

\begin{definition}[Standing Hypotheses for the Completion Results]
    \label{def:completion-standing-hypotheses}
    \lean{Physicslib4.CStarCompletion}
    \leanok
    Throughout the remainder of this subsection, $A$ denotes a type carrying
    \begin{itemize}
        \item a \texttt{NormedRing A} structure,
        \item a \texttt{StarRing A} structure,
        \item a \texttt{NormedAlgebra \ensuremath{\mathbb{C}} A} structure,
        \item a \texttt{StarModule \ensuremath{\mathbb{C}} A} structure, i.e.\ $(c \cdot a)^{*} = \bar{c} \cdot a^{*}$, and
        \item a \texttt{CStarRing A} instance, i.e.\ the C*-inequality $\|a\|\,\|a\| \le \|a^{*}a\|$.
    \end{itemize}
    We write $\widehat{A} = \texttt{UniformSpace.Completion } A$ for its completion and $\eta : A \to \widehat{A}$ for the canonical map, which has dense range by \texttt{UniformSpace.Completion.denseRange\_coe}.

    Two remarks on the shape of this hypothesis list. First, no isometry hypothesis is imposed on $\star_A$: $\|a^{*}\| = \|a\|$ follows from the C*-inequality, because \texttt{CStarRing.to\_normedStarGroup} produces the \texttt{NormedStarGroup A} instance from it. Second, \texttt{StarModule \ensuremath{\mathbb{C}} A} and \texttt{NormedAlgebra \ensuremath{\mathbb{C}} A} are listed explicitly because they are genuinely used downstream and are not consequences of the others: the \texttt{StarModule \ensuremath{\mathbb{C}} \ensuremath{\widehat{A}}} obligation of \ref{lmm:star-extends-to-completion} reduces, on the dense range of $\eta$, to \texttt{star\_smul} in $A$, which is exactly the \texttt{StarModule \ensuremath{\mathbb{C}} A} field, and \ref{lmm:completion-normed-algebra} needs an \texttt{Algebra \ensuremath{\mathbb{C}} A} to complete. Both are supplied for the colimit by \ref{lmm:quasilocal-union-normed-star-algebra}, so instantiating at the quasilocal case in \ref{lmm:quasilocal-completion-cstar} costs nothing extra.

    \emph{Why this development exists at all.} There is no C*-completion construction anywhere in Mathlib: no enveloping C*-algebra, no universal C*-algebra of a $*$-algebra, no full or reduced group C*-algebra. This was confirmed by a sweep of all 44 files of \texttt{Analysis/CStarAlgebra/}. The single C*-norm construction present there is \texttt{Unitization}, and it is not a usable template, because it builds its norm from the left regular representation --- a fundamentally different technique from completing a given C*-norm. So the results of this subsection have to be proved rather than cited. The compensation, established by the same audit and recorded node by node below, is that each of them is short: the recurring move is \texttt{UniformSpace.Completion.induction\_on} together with \texttt{isClosed\_eq} or \texttt{isClosed\_le}, pushing the claim through \texttt{norm\_coe}, \texttt{coe\_mul}, \texttt{coe\_add} and \texttt{coe\_smul} to the corresponding law on $A$.
\end{definition}

\begin{definition}[The Involution on a Completion]
    \label{def:completion-star}
    \lean{Physicslib4.instStarCompletion, Physicslib4.star_completion_coe}
    \uses{def:completion-standing-hypotheses}
    \leanok
    Let $A$ be as in \ref{def:completion-standing-hypotheses}. Define the involution on $\widehat{A}$ by
    \begin{align}
        \star_{\widehat{A}} := \texttt{UniformSpace.Completion.map } \star_{A}.
    \end{align}
    This is legitimate because $\star_A$ is an isometry --- which is not assumed but obtained from the C*-inequality on $A$ via \texttt{CStarRing.to\_normedStarGroup} --- hence uniformly continuous by \texttt{Isometry.uniformContinuous}, and \texttt{UniformSpace.Completion.map} lifts any uniformly continuous map to the completions. It is uniformly continuous, satisfies $\eta(a)^{*} = \eta(a^{*})$ by \texttt{UniformSpace.Completion.map\_coe}, and is the unique continuous map with that property by \texttt{UniformSpace.Completion.map\_unique}.

    \emph{Why this is by hand, and only here.} For the \emph{completion} the involution really does have to be written down, and Mathlib supplies nothing to start from. An audit against the local Mathlib clone found \emph{zero} occurrences of \texttt{Star}, \texttt{InvolutiveStar} or \texttt{StarRing} on \texttt{UniformSpace.Completion}: the files checked were \texttt{Topology/Algebra/UniformRing.lean}, \texttt{Topology/UniformSpace/Completion.lean}, \texttt{Topology/Algebra/GroupCompletion.lean}, \texttt{Analysis/Normed/Group/Completion.lean} and all of \texttt{Algebra/Star/}. The ring and norm structure, by contrast, is free: the instance is \texttt{UniformSpace.Completion.instNormedRing}, an auto-generated but perfectly citable name for the anonymous \texttt{instance [SeminormedRing A] : NormedRing (Completion A)} of \texttt{Analysis/Normed/Module/Completion.lean:75}. Note its hypothesis: \texttt{SeminormedRing}, with \emph{no} commutativity. That is worth holding next to the trap recorded in \ref{lmm:completion-normed-algebra}, where the analogous \texttt{NormedAlgebra} instance in the same file \emph{is} commutativity-gated and therefore does not fire. So the involution and its axioms are the whole of what this node and \ref{lmm:star-extends-to-completion} have to build. This is in contrast to the colimit, where Mathlib's \texttt{DirectLimit} does supply \texttt{Star} and \texttt{StarRing} instances --- recorded at the start of this subsection --- and nothing has to be built at all.

    \emph{The right tool, and the wrong one.} The tool to use is \texttt{UniformSpace.Completion.map}, characterised on the canonical image by \texttt{UniformSpace.Completion.map\_coe} as above. What must \emph{not} be used is \texttt{UniformSpace.Completion.mapRingHom}, even though it looks like the natural fit: it transports a \emph{ring homomorphism} to the completions, and $\star$ is anti-multiplicative, not multiplicative, so it is not a ring homomorphism $A \to A$ and does not typecheck as input. If a bundled-morphism formulation is wanted anyway, the opposite-algebra dodge is available --- read $\star$ as a ring homomorphism $A \to A^{\mathrm{op}}$ --- and it is supported, since \texttt{MulOpposite} carries a \texttt{CStarAlgebra} instance (\texttt{MulOpposite.instCStarAlgebra}). That is an optional convenience, not a requirement: the unbundled \texttt{map} route above suffices.

    \emph{Why the continuity side goals are free, recorded once for the whole block.} \texttt{UniformSpace.Completion.uniformContinuous\_map}, \texttt{UniformContinuous (Completion.map f)}, is \emph{unconditional} --- it needs no hypothesis on $f$ whatsoever, not even uniform continuity --- and it is tagged \texttt{@[fun\_prop]}, as is its corollary \texttt{UniformSpace.Completion.continuous\_map} (both in \texttt{Topology/UniformSpace/Completion.lean}, lines 480--486). Consequently every continuity side goal raised by an \texttt{isClosed\_eq} or \texttt{isClosed\_le} in this subsection --- and each of \ref{lmm:star-extends-to-completion} and \ref{lmm:completion-cstar-identity} raises at least one --- is discharged by \texttt{fun\_prop} in a single line. This is the fact that makes the four-lines-per-obligation estimate quoted throughout this block real rather than optimistic; Mathlib's own \texttt{norm\_mul\_le} field for the completion is written in exactly that style and is four lines long.
\end{definition}

\begin{lemma}[The Involution Extends to the Completion]
    \label{lmm:star-extends-to-completion}
    \lean{Physicslib4.instStarRingCompletion, Physicslib4.instStarModuleCompletion}
    \uses{def:completion-standing-hypotheses, def:completion-star}
    \leanok
    Let $A$ be as in \ref{def:completion-standing-hypotheses}. Then the operation $\star_{\widehat{A}}$ of \ref{def:completion-star} is involutive, additive, anti-multiplicative and conjugate-linear over $\mathbb{C}$,
    \begin{align}
        x^{**} = x, \quad (x+y)^{*} = x^{*} + y^{*}, \quad (xy)^{*} = y^{*}x^{*}, \quad (c \cdot x)^{*} = \bar{c} \cdot x^{*}
    \end{align}
    for $x, y \in \widehat{A}$ and $c \in \mathbb{C}$; that is, $\widehat{A}$ carries a \texttt{StarRing \ensuremath{\widehat{A}}} instance together with a \texttt{StarModule \ensuremath{\mathbb{C}} \ensuremath{\widehat{A}}} instance. The canonical map $\eta : A \to \widehat{A}$ is then a $*$-homomorphism with dense range, and the extension is the unique continuous one.

    All four laws are collected in this one node because they share a single proof skeleton --- \texttt{UniformSpace.Completion.induction\_on} plus \texttt{isClosed\_eq} plus the characterisation $\eta(a)^{*} = \eta(a^{*})$ of \ref{def:completion-star} --- and differ only in which coercion lemma and which component law of $A$ are cited at the end. That they end up bundled into two different typeclass instances, \texttt{StarRing \ensuremath{\widehat{A}}} and \texttt{StarModule \ensuremath{\mathbb{C}} \ensuremath{\widehat{A}}}, is a packaging detail and not a reason to split the mathematical content across nodes.
\end{lemma}
\begin{proof}
    \leanok
    \uses{def:completion-standing-hypotheses, def:completion-star}
    Four component laws are asserted, and they are proved as four \emph{standalone} named lemmas, not as inline field bodies: call them \texttt{Completion.star\_star'}, \texttt{Completion.star\_add'}, \texttt{Completion.star\_mul'} and \texttt{Completion.star\_smul'}. Each is about four lines by one shared skeleton. Both sides of the law are continuous in their arguments, by continuity of $\star_{\widehat{A}}$ (\ref{def:completion-star}) together with continuity of the relevant operation on the normed ring $\widehat{A}$; so the locus where they agree is closed by \texttt{isClosed\_eq} --- whose continuity side goals are the \texttt{fun\_prop} one-liners recorded in \ref{def:completion-star} --- and \texttt{UniformSpace.Completion.induction\_on} (in its \ensuremath{{}_2} form for the two-variable laws) reduces the claim to the canonical image. There the characterisation $\eta(a)^{*} = \eta(a^{*})$ of \ref{def:completion-star} (\texttt{UniformSpace.Completion.map\_coe}) applies, and the four differ only in which coercion lemma and which component law of $A$ are cited afterwards:
    \begin{itemize}
        \item \texttt{star\_star'}: \texttt{star\_star} in $A$;
        \item \texttt{star\_add'}: \texttt{UniformSpace.Completion.coe\_add} then \texttt{star\_add} in $A$;
        \item \texttt{star\_mul'}: \texttt{UniformSpace.Completion.coe\_mul} then \texttt{star\_mul} in $A$;
        \item \texttt{star\_smul'}: \texttt{UniformSpace.Completion.coe\_smul}, giving $c \cdot \eta(a) = \eta(c \cdot a)$, then \texttt{star\_smul} in $A$ --- which is exactly the \texttt{StarModule \ensuremath{\mathbb{C}} A} hypothesis of \ref{def:completion-standing-hypotheses}, and the reason that hypothesis is on the list.
    \end{itemize}
    The two instances are then assembled from those named lemmas by two one-line \texttt{where}s: \texttt{StarRing \ensuremath{\widehat{A}}} from the first three, \texttt{StarModule \ensuremath{\mathbb{C}} \ensuremath{\widehat{A}}} from the fourth. Writing the four out as lemmas rather than inlining them as field bodies is what keeps each piece at the four-line size; the assembly itself carries no content. Note that multiplication on $\widehat{A}$ is continuous but not uniformly continuous, so the anti-multiplicativity step is genuinely a continuity-plus-density argument and not a uniform-extension one; the skeleton above is chosen for that reason. Mathlib's own \texttt{norm\_mul\_le} field for \texttt{UniformSpace.Completion.instNormedRing} is written in precisely this style --- \texttt{induction\_on\ensuremath{{}_2}}, then \texttt{isClosed\_le} discharged by \texttt{fun\_prop}, then \texttt{simpa only [\ensuremath{\leftarrow} coe\_mul, norm\_coe]} --- and is four lines long, so it is the in-tree template for every obligation in this block. The ring and norm structure on $\widehat{A}$ is that same \texttt{instNormedRing} (\ref{def:completion-star}), found by typeclass inference; that $\eta$ commutes with $\star$ and has dense range, and that the extension is the unique continuous one, are recorded in \ref{def:completion-star}.

    Isometry of $\star_{\widehat{A}}$ is deliberately not part of the statement and needs no density argument of its own: once $\widehat{A}$ carries the C*-inequality (\ref{lmm:completion-cstar-identity}) the \texttt{NormedStarGroup \ensuremath{\widehat{A}}} instance, and with it $\|x^{*}\| = \|x\|$, is produced by \texttt{CStarRing.to\_normedStarGroup}.
\end{proof}

\begin{lemma}[The C*-Inequality Passes to the Completion]
    \label{lmm:completion-cstar-identity}
    \lean{Physicslib4.instCStarRingCompletion}
    \uses{def:completion-standing-hypotheses, lmm:star-extends-to-completion}
    \leanok
    Let $A$ be as in \ref{def:completion-standing-hypotheses}. Then $\widehat{A}$ satisfies
    \begin{align}
        \|x\| \, \|x\| \le \|x^{*} x\|
    \end{align}
    for all $x \in \widehat{A}$; that is, $\widehat{A}$ carries a \texttt{CStarRing} instance.

    Again only the inequality is asserted, since it is exactly the \texttt{norm\_mul\_self\_le} field of \texttt{CStarRing} and hence all there is to prove. The equality $\|x^{*}x\| = \|x\|^{2}$ follows from it by \texttt{CStarRing.norm\_star\_mul\_self}, and $\|x^{*}\| = \|x\|$ by \texttt{CStarRing.to\_normedStarGroup}.
\end{lemma}
\begin{proof}
    \leanok
    \uses{lmm:star-extends-to-completion}
    Routine, not hard, and worth saying why. \texttt{CStarRing} is a \texttt{Prop}-valued class with exactly one field, \texttt{norm\_mul\_self\_le : \ensuremath{\forall\, x,\ \|x\| * \|x\| \le \|x^{\star} * x\|}} (\texttt{Analysis/CStarAlgebra/Basic.lean:89}). That field is a \emph{non-strict inequality} between two continuous real-valued functions of the single variable $x$, so it transports to the completion in about four lines: the locus where it holds is closed by \texttt{isClosed\_le}, and \texttt{UniformSpace.Completion.induction\_on} reduces the goal to the canonical image.

    On the canonical image it is the C*-inequality of $A$: $\eta$ is a ring and $*$-homomorphism (\ref{lmm:star-extends-to-completion}) and is norm preserving by \texttt{UniformSpace.Completion.norm\_coe}, so $\|\eta(a)^{*}\eta(a)\| = \|\eta(a^{*}a)\| = \|a^{*}a\| \ge \|a\|\,\|a\| = \|\eta(a)\|\,\|\eta(a)\|$. Continuity of the two sides is the same package as elsewhere: the involution is continuous by \ref{lmm:star-extends-to-completion}, and multiplication and the norm are continuous on a normed ring.

    Two consequences worth recording. First, isometry of the involution on $\widehat{A}$ is \emph{not} a separate obligation anywhere in this block: it follows from this inequality through \texttt{CStarRing.to\_normedStarGroup} (\texttt{Basic.lean:114}). Second, \texttt{CStarRing} does not itself require completeness --- its hypotheses are only \texttt{NonUnitalNormedRing} and \texttt{StarRing} --- which is why the same inequality is available on the colimit before completing, as \ref{lmm:quasilocal-colimit-cstar-identity} asserts.
\end{proof}

\begin{lemma}[The Completion is a Normed $\mathbb{C}$-Algebra]
    \label{lmm:completion-normed-algebra}
    \lean{Physicslib4.instNormedAlgebraCompletion}
    \uses{def:completion-standing-hypotheses}
    \leanok
    Let $A$ be as in \ref{def:completion-standing-hypotheses}. Then $\widehat{A}$ carries a \texttt{NormedAlgebra \ensuremath{\mathbb{C}} \ensuremath{\widehat{A}}} instance.
\end{lemma}
\begin{proof}
    \leanok
    \uses{def:completion-standing-hypotheses}
    This node exists to work around one specific trap, and the trap is the whole of its content. Mathlib \emph{does} have a \texttt{NormedAlgebra} instance on a completion (\texttt{Analysis/Normed/Module/Completion.lean:87}), but it is gated on \texttt{SeminormedCommRing A} --- \emph{commutative} --- so it does \emph{not} fire here: $A$ is noncommutative in every case of interest. Without noticing this one would expect the instance for free and discover at formalization time that typeclass inference simply fails to find it.

    Supplying it by hand is short. \texttt{NormedAlgebra} has exactly one field beyond \texttt{Algebra}, namely \texttt{norm\_smul\_le}, $\|c \cdot x\| \le |c|\,\|x\|$; and that field is already available from the \texttt{NormedSpace} instance on the completion of a normed space (line 36 of the same file), which carries no commutativity hypothesis. So the instance is \texttt{UniformSpace.Completion.algebra} together with \texttt{norm\_smul\_le} borrowed from that \texttt{NormedSpace} instance --- a \texttt{where} with a single field.

    This is also where scalar-action compatibility would live, if it is needed at all, and the evidence is that it is not. There are two a priori different \texttt{SMul \ensuremath{\mathbb{C}} \ensuremath{\widehat{A}}} structures in play: the one underlying \texttt{UniformSpace.Completion.algebra}, characterised by its \texttt{smul\_def'} field as $c \cdot x = \eta(\texttt{algebraMap } c)\, x$, and the one underlying the \texttt{NormedSpace} instance, which extends the scalar action on $A$. Mathlib's own commutative \texttt{NormedAlgebra} instance has the body \texttt{norm\_smul\_le := norm\_smul\_le} and nothing else: it borrows the \texttt{NormedSpace} field directly, with \emph{no} agreement argument, and --- decisively --- \emph{never uses the commutativity hypothesis in that body}. So the identical body compiles in the noncommutative setting, and this node is three lines. Nothing below is needed; it is recorded only against the remote possibility that the two actions fail to be definitionally equal, in which case the fix is the usual \texttt{induction\_on} plus \texttt{isClosed\_eq} step: on the canonical image each action sends $\eta(a)$ to $\eta(c \cdot a)$, the first by \texttt{UniformSpace.Completion.coe\_mul} with \texttt{Algebra.smul\_def} in $A$, the second by \texttt{UniformSpace.Completion.coe\_smul}.

    For the record, so that none of it is re-proved: what \emph{is} free on the completion, without commutativity, is \texttt{UniformSpace.Completion.ring} and \texttt{UniformSpace.Completion.algebra} --- both named declarations, in \texttt{Topology/Algebra/UniformRing.lean} --- together with \texttt{UniformSpace.Completion.instNormedRing}, the \texttt{NormedSpace} and \texttt{NormedAddCommGroup} instances of \texttt{Analysis/Normed/Module/Completion.lean}, and the norm-preservation lemma \texttt{UniformSpace.Completion.norm\_coe}. Those three instances are declared anonymously in the source but carry the usual auto-generated names and can be cited by them; and none is commutativity-gated --- \texttt{instNormedRing} asks only for \texttt{SeminormedRing A}, and the \texttt{NormedSpace} instance only for a seminormed additive group with a scalar action. It is precisely and only the \texttt{NormedAlgebra} instance twelve lines further down the same file that adds \texttt{SeminormedCommRing}, which is what makes this node necessary.
\end{proof}

\begin{lemma}[The Completion of a C*-Normed $*$-Algebra is a C*-Algebra]
    \label{lmm:completion-of-cstar-normed-star-algebra}
    \lean{Physicslib4.instCStarAlgebraCompletion}
    \uses{def:completion-standing-hypotheses, lmm:star-extends-to-completion, lmm:completion-cstar-identity, lmm:completion-normed-algebra}
    \leanok
    Let $A$ be as in \ref{def:completion-standing-hypotheses}. Then its completion $\widehat{A}$ is a C*-algebra.
\end{lemma}
\begin{proof}
    \leanok
    \uses{lmm:star-extends-to-completion, lmm:completion-cstar-identity, lmm:completion-normed-algebra}
    Pure bundling, and the audit confirms it is literally that: \texttt{CStarAlgebra} is declared as \texttt{class CStarAlgebra (A : Type*) extends NormedRing A, StarRing A, CompleteSpace A, CStarRing A, NormedAlgebra \ensuremath{\mathbb{C}} A, StarModule \ensuremath{\mathbb{C}} A} with an \emph{empty} body (\texttt{Analysis/CStarAlgebra/Classes.lean:30-42}). It adds no field of its own beyond its parents, so once the parents below are in place this node is a \texttt{where} with an essentially empty body. Exactly those six parents are supplied for $\widehat{A}$:
    \begin{itemize}
        \item \texttt{NormedRing \ensuremath{\widehat{A}}}: found by typeclass inference, and citable by name as \texttt{UniformSpace.Completion.instNormedRing} (\texttt{Analysis/Normed/Module/Completion.lean:75}), whose only hypothesis is \texttt{SeminormedRing A} --- no commutativity, unlike the \texttt{NormedAlgebra} instance in the same file that forces \ref{lmm:completion-normed-algebra} to exist. Note that multiplication on $A$ is \emph{not} uniformly continuous --- only uniformly continuous on bounded sets --- so it is extended as a bounded bilinear map rather than by uniform continuity, which is exactly what that instance arranges.
        \item \texttt{StarRing \ensuremath{\widehat{A}}}: \ref{lmm:star-extends-to-completion}.
        \item \texttt{CompleteSpace \ensuremath{\widehat{A}}}: nothing to prove --- this is the ambient \texttt{CompleteSpace} instance that comes with the completion, \texttt{UniformSpace.Completion.completeSpace}, not an obligation of this node.
        \item \texttt{CStarRing \ensuremath{\widehat{A}}}: \ref{lmm:completion-cstar-identity}.
        \item \texttt{NormedAlgebra \ensuremath{\mathbb{C}} \ensuremath{\widehat{A}}}: \ref{lmm:completion-normed-algebra}.
        \item \texttt{StarModule \ensuremath{\mathbb{C}} \ensuremath{\widehat{A}}}: \ref{lmm:star-extends-to-completion}, which supplies this instance alongside the \texttt{StarRing} one above --- the conjugate-linearity law is one of the four it establishes.
    \end{itemize}
    The \texttt{StarModule} field was missing from an earlier version of this node; it is listed here so that the obligation is not discovered only at formalization time.
\end{proof}

\begin{lemma}[The Completion of the Quasilocal Colimit is a C*-Algebra]
    \label{lmm:quasilocal-completion-cstar}
    \lean{Physicslib4.AQFT.HaagKastler.QuasilocalCompletion, Physicslib4.AQFT.HaagKastler.quasilocalCompletionCStarAlgebra}
    \uses{lmm:quasilocal-union-normed-star-algebra, lmm:quasilocal-colimit-cstar-identity, def:completion-standing-hypotheses, lmm:completion-of-cstar-normed-star-algebra}
    \leanok
    The completion of the normed $*$-algebra of \ref{lmm:quasilocal-union-normed-star-algebra} is a C*-algebra.
\end{lemma}
\begin{proof}
    \leanok
    \uses{lmm:quasilocal-union-normed-star-algebra, lmm:quasilocal-colimit-cstar-identity, def:completion-standing-hypotheses, lmm:completion-of-cstar-normed-star-algebra}
    This node is where the general completion results are instantiated at the colimit. Take $A$ to be the normed $*$-algebra of \ref{lmm:quasilocal-union-normed-star-algebra}. Its \texttt{NormedRing}, \texttt{StarRing}, \texttt{NormedAlgebra \ensuremath{\mathbb{C}}} and \texttt{StarModule \ensuremath{\mathbb{C}}} structures are exactly what that lemma supplies, and the remaining hypothesis of \ref{def:completion-standing-hypotheses}, the C*-inequality, is \ref{lmm:quasilocal-colimit-cstar-identity}. So \ref{lmm:completion-of-cstar-normed-star-algebra} applies to $A$ verbatim, and the claim is that lemma applied to it. In particular isometry of the involution is not a separate assumption to discharge: it comes from the same inequality through \texttt{CStarRing.to\_normedStarGroup}.
\end{proof}

Two matters of project bookkeeping attach to \ref{lmm:quasilocal-completion-cstar} but are not part of its proof, and are therefore recorded here as surrounding prose rather than inside it.

\emph{Uniqueness of the complete C*-norm: deliberately not claimed.} An earlier version of that node also asserted that the resulting C*-norm is the only complete C*-norm on the carrier, citing Mathlib's \texttt{StarAlgEquiv.norm\_map}, \texttt{lemma norm\_map (\ensuremath{\varphi} : F) (a : A) : \ensuremath{\|\varphi a\| = \|a\|}}. That citation does not support the claim. \texttt{StarAlgEquiv.norm\_map} is about a $*$-isomorphism $\varphi$ between \emph{two} types, each already carrying its own C*-algebra structure; it says such a map is isometric, with no injectivity hypothesis needed. Uniqueness of the complete C*-norm is a statement about \emph{one} carrier equipped with two \texttt{CStarAlgebra} structures whose underlying ring, star and $\mathbb{C}$-algebra structures agree, concluding that the two \texttt{Norm} fields are equal. Getting from the former to the latter needs the identity map to be exhibited as a \texttt{StarAlgEquiv} between two such structures, which in Lean means type synonyms and instance plumbing --- real work, not a rewrite. The clause is therefore dropped from the statement of that node and recorded here as prose rather than left as an unbacked assertion; it is available to be added later as a declaration of its own, with the two-structures-on-one-carrier statement written out, if a consumer needs it.

The completeness caveat is worth keeping. The completeness hypothesis is hidden in the word ``C*-algebra'' and is essential: a $*$-algebra can carry more than one C*-norm, so what could be unique is only the \emph{complete} one. The counterexample is the group $*$-algebra $\mathbb{C}[F_2]$ of the free group on two generators, which carries the distinct full and reduced C*-norms, neither of them complete. An earlier version of this blueprint claimed that any two C*-norms on a $*$-algebra coincide, which is that false statement.

This is the Bratteli--Robinson 2.2.6 material that Chapter 6.2 discusses informally. Uniqueness is not a mere technicality there: it is what would make the quasilocal algebra well defined \emph{up to isomorphism} rather than merely existing, since without it the completion could depend on a choice of norm and different choices could give non-isomorphic C*-algebras. What \ref{lmm:quasilocal-completion-cstar} establishes is existence; the well-definedness reading rests on the prose above.

\emph{The construction route is settled: colimit, then completion.} The project constructs the quasilocal algebra \emph{from the net alone} --- the directed colimit of the local C*-algebras along the isotony family, followed by completion --- which is exactly the chain running from \ref{def:completion-standing-hypotheses} through \ref{lmm:quasilocal-completion-cstar}. This is a decision that has been taken, not an option still under consideration, and the reason is physical. The alternative below is rejected.

\emph{The rejected alternative, and why it is rejected.} The alternative was to realise the quasilocal algebra as a \emph{closed} $*$-subalgebra of an ambient C*-algebra: take the $*$-subalgebra generated by the images of the canonical embeddings $\iota_{\mathbf{B}}$ --- data that Axiom 3 (\ref{def:local-commutativity}) already assumes --- and close it off with \texttt{StarSubalgebra.topologicalClosure}. That description is accurate and the saving would have been real: \texttt{StarSubalgebra.cstarAlgebra} supplies the \emph{entire} C*-structure as an instance --- involution, norm, completeness, the C*-identity, the algebra structure and the \texttt{StarModule} field all at once --- so \ref{def:completion-standing-hypotheses}, \ref{def:completion-star}, \ref{lmm:star-extends-to-completion}, \ref{lmm:completion-cstar-identity}, \ref{lmm:completion-normed-algebra} and \ref{lmm:completion-of-cstar-normed-star-algebra} would all have become unnecessary, with none of the density arguments above written by hand.

It is rejected because of what it presupposes. \texttt{StarSubalgebra.cstarAlgebra} takes an \emph{ambient} \texttt{CStarAlgebra A} as a hypothesis and equips a closed $*$-subalgebra of it with the induced structure; it does not produce a C*-algebra out of nothing. So the shortcut requires an ambient C*-algebra containing isomorphic copies of every local algebra $\mathfrak{U}(\mathbf{B})$ to be given in advance --- and there is no physical justification for such an algebra. Nothing in the physics supplies one, nothing says where it would come from, so it cannot be assumed. That is the whole of the argument: the shortcut would have the blueprint assume $\mathfrak{U}$ rather than construct it from the net, and assuming it has no basis in the physics. The completion chain above starts from the net alone and \emph{builds} a C*-algebra, discharging the existence claim that \ref{def:quasilocal-algebra} makes.

\emph{The existing route-2 material in this repository is not invalidated.} \texttt{Physicslib4.AQFT.HaagKastler.dense\_adjoin\_iUnion\_range\_\ensuremath{\iota}} (in \texttt{QuasilocalIntertwiner.lean}) and \texttt{Physicslib4.exists\_starAlgHom\_extend\_of\_dense} (in \texttt{Physicslib4/Analysis/CStarDenseExtend.lean}) implement the dense-extension technique and both take the quasilocal algebra as given. Taking it as given is perfectly legitimate for a \emph{consumer} of the construction: these are downstream results about a quasilocal algebra, not competing constructions of one, and they remain valid verbatim. What the decision changes is only that the blueprint must now \emph{also} construct such an algebra, which is what the chain above does. Indeed \texttt{CStarDenseExtend.lean} remains the closest worked example in this repository of transporting star-structure across a dense embedding, so it is the reference to read before formalizing the completion nodes above.

Before introducing the next axiom, we must introduce the definition:

\begin{definition}[Quasilocal Algebra]
    \label{def:quasilocal-algebra}
    \lean{Physicslib4.AQFT.HaagKastler.QuasilocalAlgebra}
    \leanfile{Physicslib4/AQFT/HaagKastler/QuasilocalAlgebra.lean}
    \leanok
    \uses{def:local-algebras, def:isotony, lmm:minkowski-diamonds-upward-directed, lmm:quasilocal-union-normed-star-algebra, lmm:quasilocal-completion-cstar}
    Consider the union of all $\mathfrak{U}(\mathbf{B})$, taken along the isotony family of Axiom 2 (\ref{def:isotony}). This union is a normed *-algebra; taking its completion one obtains a C*-algebra denoted $\mathfrak{U}$, called the \textit{quasilocal algebra}.

    The union is to be read as a \emph{directed colimit}, not as a set-theoretic union. This is what Axiom 2's identity and composition laws buy: with $\mathbf{B} \mapsto \mathfrak{U}(\mathbf{B})$ a functor on the inclusion order, and the Alexandrov diamonds upward directed, the local algebras form a directed system and the colimit carries a well-defined $*$-algebra structure --- the product of $a_1 \in \mathfrak{U}(\mathbf{B}_1)$ and $a_2 \in \mathfrak{U}(\mathbf{B}_2)$ is computed in any $\mathfrak{U}(\mathbf{B})$ containing both, and functoriality is exactly what makes the answer independent of that choice. Without the two laws there is no such structure: a bare set-theoretic union of the $\mathfrak{U}(\mathbf{B})$ has no multiplication at all, since elements of different local algebras live in unrelated carriers.

    The three supporting results this definition rests on are now declarations of this blueprint rather than appeals to the informal discussion of Chapters 6.1 and 6.2: (i) upward directedness of the Alexandrov diamonds (\ref{lmm:minkowski-diamonds-upward-directed}), which supplies the containing diamond $\mathbf{B}$ in the description above; (ii) that the colimit of the directed system is a normed $*$-algebra (\ref{lmm:quasilocal-union-normed-star-algebra}); and (iii) that its completion is a C*-algebra (\ref{lmm:quasilocal-completion-cstar}). Of these, (i) is already formalized, while (ii) and (iii) remain open. What would pin $\mathfrak{U}$ down up to isomorphism rather than merely produce it is uniqueness of the complete C*-norm; that is discussed as prose in the passage following \ref{lmm:quasilocal-completion-cstar} and is deliberately not claimed by any declaration here, for the reasons given there.

    That this definition is marked as formalized while (ii) and (iii) are not is not an inconsistency: the Lean takes $\mathfrak{U}$ as \emph{given} --- an ambient C*-algebra together with the canonical embeddings $\iota_{\mathbf{B}}$ and their cocone condition, as recorded in \ref{def:local-commutativity} --- rather than constructing it as a colimit and completing it. Lemmas (ii) and (iii) are the obligations one incurs by \emph{building} such a $\mathfrak{U}$ from the local algebras; they are not gaps in the existing Lean. See the formalization note following \ref{lmm:quasilocal-completion-cstar} on the two available construction routes.
\end{definition}

\begin{definition}[Axiom 3: Local Commutativity]
    \label{def:local-commutativity}
    \lean{Physicslib4.AQFT.HaagKastler.LocalCommutativity}
    \leanfile{Physicslib4/AQFT/HaagKastler/LocalCommutativity.lean}
    \leanok
    \uses{def:alexandrov-topology, def:minkowski-spacetime, def:chronological-future-and-chronological-past, def:local-algebras, def:isotony, def:completely-spacelike, def:quasilocal-algebra}
    Let $\mathbf{B}_1$ and $\mathbf{B}_2$ be any two basis elements of the Alexandrov topology on Minkowski spacetime, i.e. any two sets of the form $I^+(p_1) \cap I^-(q_1)$ and $I^+(p_2) \cap I^-(q_2)$.

    If $\mathbf{B_1}$ and $\mathbf{B_2}$ are completely spacelike, then $\mathfrak{U}(\mathbf{B_1})$ and $\mathfrak{U}(\mathbf{B_2})$ commute in the quasilocal algebra $\mathfrak{U}$, i.e. for any $a_1$ in $\mathfrak{U}(\mathbf{B_1})$ and $a_2$ in $\mathfrak{U}(\mathbf{B_2})$ it follows that
    \begin{align}
        \iota_{\mathbf{B}_1}(a_1)\, \iota_{\mathbf{B}_2}(a_2) - \iota_{\mathbf{B}_2}(a_2)\, \iota_{\mathbf{B}_1}(a_1) = 0
    \end{align}
    in the quasilocal algebra $\mathfrak{U}$. Here $\iota_{\mathbf{B}} : \mathfrak{U}(\mathbf{B}) \to \mathfrak{U}$ is the \emph{canonical embedding} of a local algebra into the quasilocal algebra of \ref{def:quasilocal-algebra} --- the map into the completion of the union of the local algebras. It is \emph{not} the unital $*$-monomorphism $i_{\mathbf{B}_1\mathbf{B}_2} : \mathfrak{U}(\mathbf{B}_1) \hookrightarrow \mathfrak{U}(\mathbf{B}_2)$ of Axiom 2 (Isotony), which goes from one local algebra to another and not into $\mathfrak{U}$.

    The two families are required to be compatible: for all basis sets $\mathbf{B}_1 \subseteq \mathbf{B}_2$,
    \begin{align}
        \iota_{\mathbf{B}_2} \circ i_{\mathbf{B}_1\mathbf{B}_2} = \iota_{\mathbf{B}_1}.
    \end{align}
    This is the cocone condition making the $\iota_{\mathbf{B}}$ a compatible family on the directed system of Axiom 2 (\ref{def:isotony}), and it is what makes $\iota_{\mathbf{B}}$ well defined on the colimit: an element of $\mathfrak{U}(\mathbf{B}_1)$ may be regarded as an element of any larger $\mathfrak{U}(\mathbf{B}_2)$, and all such readings must have the same image in $\mathfrak{U}$. Without it the displayed commutator would depend on which local algebra $a_1$ and $a_2$ were viewed in.

    Note that the curved counterpart (\ref{def:local-commutativity-in-curved-spacetime}) has a different shape: there is no quasilocal algebra in curved spacetime, so commutation is stated inside a common containing algebra $\mathfrak{U}(\mathbf{B})$ using the Axiom 2 embeddings directly, and no $\iota$ appears.
\end{definition}

The next axiom makes use of the following new definition:

\begin{definition}[Quasilocal Observable]
    \label{def:quasilocal-observable}
    \lean{Physicslib4.AQFT.HaagKastler.IsQuasilocalObservable}
    \leanfile{Physicslib4/AQFT/HaagKastler/QuasilocalCompleteness.lean}
    \leanok
    \uses{def:quasilocal-algebra, thrm:gns-construction-theorem, def:state}
    The image $\pi_\omega(a)$ of a self-adjoint member $a$ of the quasilocal algebra $\mathfrak{U}$ under a GNS *-homomorphism $\pi_\omega$ is self-adjoint and thus corresponds to an ``observable''. Any ``observable'' corresponding to such a self-adjoint $\pi_\omega(a)$ is called a \textit{quasilocal observable}.
\end{definition}

\begin{definition}[Axiom 4: Quasilocal Completeness]
    \label{def:quasilocal-completeness}
    \lean{Physicslib4.AQFT.HaagKastler.QuasilocalCompleteness}
    \leanfile{Physicslib4/AQFT/HaagKastler/QuasilocalCompleteness.lean}
    \leanok
    \uses{def:quasilocal-observable}
    All ``observables'' are quasilocal observables.
\end{definition}

\begin{definition}[Axiom 5: Lorentz Covariance]
    \label{def:lorentz-covariance}
    \lean{Physicslib4.AQFT.HaagKastler.LorentzCovariance}
    \leanfile{Physicslib4/AQFT/HaagKastler/LorentzCovariance.lean}
    \leanok
    \uses{def:alexandrov-topology, def:minkowski-spacetime, def:chronological-future-and-chronological-past, def:local-algebras, def:isotony}
    Let $\mathbf{B}$ be any basis element of the Alexandrov topology on Minkowski spacetime, i.e. any set of the form $I^+(p) \cap I^-(q)$.

    A member $L$ of the inhomogeneous Lorentz group connected to the identity acts on $\mathfrak{U}(\mathbf{B})$ as follows
    \begin{align}
        \alpha_L : \mathfrak{U}(\mathbf{B}) &\longrightarrow \mathfrak{U}(L\mathbf{B}) \\
                             a              &\longmapsto     \alpha_L(a)
    \end{align}
    where $L\mathbf{B}$ is the image of the region $\mathbf{B}$ under the transformation $L$ and $\alpha_L$ is a unital *-isomorphism generated by $L$. The map $\alpha_L$ is such that (1) for the identity element $\mathbf{1}$ of the Lorentz group it satisfies
    \begin{align}
        \alpha_{\mathbf{1}}(a) = a,
    \end{align}
    (2) for all appropriate $a$, $L$, and $L'$ it satisfies
    \begin{align}
        \alpha_{L' \cdot L}(a) = \alpha_{L'}(\alpha_{L}(a)),
    \end{align}
    and (3) for basis elements $\mathbf{B}_\iota \subset \mathbf{B}_\kappa$ and the unital *-monomorphism $i$ of Axiom 2 (Isotony) $\alpha_L$ commutes with $i$. In other words the following diagram

    \begin{center}
        \begin{tikzcd}
            \mathfrak{U}(\mathbf{B}_\kappa) \arrow[r, "\alpha_L"]
            & \mathfrak{U}(L\mathbf{B}_\kappa) \\
            \mathfrak{U}(\mathbf{B}_\iota) \arrow[r, "\alpha_L"]  \arrow[u, "i"]
            & \mathfrak{U}(L\mathbf{B}_\iota)  \arrow[u, "i"]
        \end{tikzcd}
    \end{center}

    commutes.
\end{definition}

\begin{definition}[Haag-Kastler Net]
    \label{def:haag-kastler-net}
    \lean{Physicslib4.AQFT.HaagKastler.HaagKastlerNet}
    \leanfile{Physicslib4/AQFT/HaagKastler/Net.lean}
    \uses{def:local-algebras, def:isotony, def:local-commutativity, def:quasilocal-completeness, def:lorentz-covariance}
    \leanok
    A \emph{Haag-Kastler net} on Minkowski spacetime is the bundling of the data of \ref{def:local-algebras} together with the properties of \ref{def:isotony}, \ref{def:local-commutativity}, \ref{def:quasilocal-completeness}, and \ref{def:lorentz-covariance}. In the Lean formalization this is a single structure whose fields are the assignment $\mathbf{B} \mapsto \mathfrak{U}(\mathbf{B})$ and proofs that this assignment satisfies the four remaining axioms. Theorems about AQFT take an instance of this structure as a hypothesis and invoke each axiom as a projection.
\end{definition}

This concludes the presentation of the ``sharpened'' axioms. As proven in this blog post, these ``sharpened'' axioms entail the original axioms save Axiom 6 (Primitivity), which is an axiom that we abandon.

These ``sharpened'' axioms also allow for a straightforward generalization to a set of axioms describing AQFT in curved spacetime, i.e. a set of axioms describing AQFT in a curved spacetime background that is fixed and treated classically. In a subsequent blog post we will detail these axioms.

\subsection{Einstein Causality}

Local commutativity (\ref{def:local-commutativity}) is an algebraic statement about the quasilocal algebra. Its physical content - that spacelike-separated measurements do not interfere - is its operator form: in any representation of the quasilocal algebra, spacelike-separated local observables commute as bounded operators on the Hilbert space.

\begin{theorem}[Einstein Causality in a Representation]
    \label{thrm:einstein-causality}
    \lean{Physicslib4.AQFT.HaagKastler.HaagKastlerNet.einstein_causality, Physicslib4.AQFT.HaagKastler.HaagKastlerNet.exists_gns_einstein_causality}
    \leanfile{Physicslib4/AQFT/HaagKastler/EinsteinCausality.lean}
    \uses{def:haag-kastler-net, thrm:gns-construction-theorem, def:state}
    \leanok
    Let $\pi$ be any $*$-representation of the quasilocal algebra on a Hilbert space $H$. If $\mathbf{B}_1, \mathbf{B}_2$ are completely spacelike-separated basis regions, then for all $a \in \mathfrak{U}(\mathbf{B}_1)$ and $b \in \mathfrak{U}(\mathbf{B}_2)$ the operators $\pi(\iota_{\mathbf{B}_1} a)$ and $\pi(\iota_{\mathbf{B}_2} b)$ commute.
\end{theorem}
\begin{proof}
    \leanok
    \uses{def:local-commutativity}
    This is the image under $\pi$ of the local commutativity relation (\ref{def:local-commutativity}); since the GNS representation of any state is such a $\pi$, spacelike-separated local observables commute on every GNS Hilbert space.
\end{proof}

\subsection{Local von Neumann Algebras}

The abstract C*-algebras $\mathfrak{U}(\mathbf{B})$ are a stepping stone; the genuine object of algebraic quantum field theory is the net of von Neumann algebras they generate in a representation. Given a $*$-representation $\pi$ of the quasilocal algebra on a Hilbert space $H$, the \emph{local von Neumann algebra} of a region $\mathbf{B}$ is the bicommutant $R(\mathbf{B}) = \pi(\mathfrak{U}(\mathbf{B}))''$ of the local observable operators. We model the commutant by the centralizer in $\mathcal{B}(H)$.

\begin{definition}[Local von Neumann Algebra]
    \label{def:local-von-neumann}
    \lean{Physicslib4.AQFT.HaagKastler.HaagKastlerNet.localOperators, Physicslib4.AQFT.HaagKastler.HaagKastlerNet.localVonNeumann}
    \leanfile{Physicslib4/AQFT/HaagKastler/LocalVonNeumann.lean}
    \leanok
    \uses{def:haag-kastler-net, thrm:gns-construction-theorem}
    Let $\pi$ be a $*$-representation of the quasilocal algebra on $H$. The \emph{local observable operators} of a region $\mathbf{B}$ are the image $\pi(\mathfrak{U}(\mathbf{B})) = \{\pi(\iota_{\mathbf{B}} a) : a \in \mathfrak{U}(\mathbf{B})\}$. The \emph{local von Neumann algebra} $R(\mathbf{B})$ is the bicommutant $\pi(\mathfrak{U}(\mathbf{B}))''$.
\end{definition}

\begin{lemma}[The Bicommutant of a Self-Adjoint Set is a von Neumann Algebra]
    \label{lmm:bicommutant-of-selfadjoint-is-von-neumann}
    \lean{Physicslib4.GNS.vonNeumannOfSelfAdjoint}
    \leanfile{Physicslib4/GNS/Irreducibility.lean}
    \leanok
    Let $S$ be a self-adjoint set of bounded operators on a Hilbert space $H$, that is, $x \in S$ implies $x^* \in S$. Then the bicommutant $S''$ is a von Neumann algebra: it is a $*$-subalgebra of $\mathcal{B}(H)$, and it is bicommutant-closed, $S'''' = S''$. This is the general construction that bundles every local algebra below, applied to the self-adjoint set of local observable operators.
\end{lemma}
\begin{proof}
    \leanok
    The commutant of a self-adjoint set is $*$-closed, which is Mathlib's \texttt{Set.star\_mem\_centralizer}: if $T \in S'$ and $x \in S$, then $x^* \in S$, so $T x^* = x^* T$, and taking adjoints gives $x T^* = T^* x$; hence $T^* \in S'$. So $S'$ is itself $*$-closed, and the centralizer of any set is a subalgebra of $\mathcal{B}(H)$ (Mathlib's \texttt{Subalgebra.centralizer}); combining the two makes $S'' = (S')'$ a $*$-subalgebra. (These two are exactly what $\mathtt{Physicslib4.GNS.starSubalgebraCentralizer}$ is assembled from. Mathlib's \texttt{StarSubalgebra.centralizer} is \emph{not} the right citation: its carrier is the centralizer of the star-closure $s \cup s^*$, not of $s$.) For bicommutant-closedness, the triple-commutant collapse $S''' = S'$ (Mathlib's \texttt{Set.centralizer\_centralizer\_centralizer}) gives, on taking one further commutant, $S'''' = (S''')' = (S')' = S''$.
\end{proof}

\begin{definition}[$R(\mathbf{B})$ as a von Neumann Algebra]
    \label{def:local-von-neumann-algebra}
    \lean{Physicslib4.AQFT.HaagKastler.HaagKastlerNet.localVonNeumannAlgebra}
    \leanfile{Physicslib4/AQFT/HaagKastler/LocalVonNeumann.lean}
    \leanok
    \uses{def:local-von-neumann, lmm:bicommutant-of-selfadjoint-is-von-neumann}
    The local algebra $R(\mathbf{B})$ is registered as a genuine \texttt{VonNeumannAlgebra} (Mathlib's bundled structure), not merely a set of operators: it is the bundled algebra supplied by \ref{lmm:bicommutant-of-selfadjoint-is-von-neumann} for the self-adjoint set $S = \pi(\mathfrak{U}(\mathbf{B}))$ of local observable operators (self-adjoint because $\pi$ and $\iota_{\mathbf{B}}$ are $*$-homomorphisms). Its underlying set is the bicommutant $\pi(\mathfrak{U}(\mathbf{B}))''$ of \ref{def:local-von-neumann}.
\end{definition}

\begin{theorem}[Microcausality at the von Neumann Level]
    \label{thrm:von-neumann-microcausality}
    \lean{Physicslib4.AQFT.HaagKastler.HaagKastlerNet.localVonNeumann_subset_centralizer}
    \leanfile{Physicslib4/AQFT/HaagKastler/LocalVonNeumann.lean}
    \uses{def:local-von-neumann}
    \leanok
    For completely spacelike-separated basis regions $\mathbf{B}_1, \mathbf{B}_2$, the local von Neumann algebras commute: $R(\mathbf{B}_1) \subseteq R(\mathbf{B}_2)'$.
\end{theorem}
\begin{proof}
    \leanok
    \uses{def:local-von-neumann, thrm:einstein-causality}
    This is the von Neumann form of Einstein causality (\ref{thrm:einstein-causality}): elementwise commutation $\pi(\mathfrak{U}(\mathbf{B}_2)) \subseteq \pi(\mathfrak{U}(\mathbf{B}_1))'$, with the antitonicity of the commutant and the identity $S''' = S'$ collapsing the iterated commutants. A Haag-Kastler net is thus a net of mutually commuting von Neumann algebras for spacelike regions.
\end{proof}

\begin{theorem}[Isotony of the von Neumann Net]
    \label{thrm:von-neumann-isotony}
    \lean{Physicslib4.AQFT.HaagKastler.HaagKastlerNet.localVonNeumann_mono}
    \leanfile{Physicslib4/AQFT/HaagKastler/LocalVonNeumann.lean}
    \uses{def:local-von-neumann}
    \leanok
    For basis regions $\mathbf{B}_1 \subseteq \mathbf{B}_2$, the local von Neumann algebras are nested: $R(\mathbf{B}_1) \subseteq R(\mathbf{B}_2)$.
\end{theorem}
\begin{proof}
    \leanok
    \uses{def:local-von-neumann, def:isotony}
    The local observables of $\mathbf{B}_1$ embed into those of $\mathbf{B}_2$ via the quasilocal isotony coherence, and the double commutant is monotone.
\end{proof}

\begin{theorem}[Bundled von Neumann Microcausality and Isotony]
    \label{thrm:von-neumann-bundled-order}
    \lean{Physicslib4.AQFT.HaagKastler.HaagKastlerNet.localVonNeumannAlgebra_le_commutant, Physicslib4.AQFT.HaagKastler.HaagKastlerNet.localVonNeumannAlgebra_mono}
    \leanfile{Physicslib4/AQFT/HaagKastler/LocalVonNeumann.lean}
    \uses{def:local-von-neumann-algebra}
    \leanok
    Phrased on the bundled \texttt{VonNeumannAlgebra} $R(\mathbf{B})$ with its order $\le$ and Mathlib's commutant: for completely spacelike-separated regions $R(\mathbf{B}_1) \le R(\mathbf{B}_2)'$, and for $\mathbf{B}_1 \subseteq \mathbf{B}_2$, $R(\mathbf{B}_1) \le R(\mathbf{B}_2)$.
\end{theorem}
\begin{proof}
    \leanok
    \uses{def:local-von-neumann-algebra, thrm:von-neumann-microcausality, thrm:von-neumann-isotony}
    Both reduce to the set-level statements through the coercion $\uparrow R(\mathbf{B}) = \pi(\mathfrak{U}(\mathbf{B}))''$.
\end{proof}

\begin{definition}[The Net of von Neumann Algebras]
    \label{def:von-neumann-net}
    \lean{Physicslib4.AQFT.HaagKastler.HaagKastlerNet.vonNeumannNet}
    \leanfile{Physicslib4/AQFT/HaagKastler/LocalVonNeumann.lean}
    \leanok
    \uses{def:local-von-neumann-algebra, thrm:von-neumann-bundled-order}
    Packaging isotony, the assignment $\mathbf{B} \mapsto R(\mathbf{B})$ is an order-preserving map from the poset of basis regions (ordered by inclusion) to the von Neumann algebras of $\mathcal{B}(H)$. This realizes \emph{the net} as a functor on the inclusion poset: containment of regions is sent to containment of algebras. It is the central object of the algebraic approach.
\end{definition}

\begin{theorem}[Statistical Independence (Schlieder Property)]
    \label{thrm:statistical-independence}
    \lean{Physicslib4.AQFT.HaagKastler.HaagKastlerNet.localVonNeumann_separating, Physicslib4.AQFT.HaagKastler.HaagKastlerNet.eq_zero_of_commute_of_cyclic}
    \leanfile{Physicslib4/AQFT/HaagKastler/LocalVonNeumann.lean}
    \uses{def:cyclic-vector}
    \leanok
    If $\Omega$ is cyclic for the local observables of $\mathbf{B}_1$, then for a spacelike-separated region $\mathbf{B}_2$ every $R \in R(\mathbf{B}_2)$ with $R\Omega = 0$ is zero, so $\Omega$ is separating for $R(\mathbf{B}_2)$. In Minkowski spacetime the cyclicity hypothesis is exactly the content of the Reeh-Schlieder theorem, which rests on the spectrum condition; here it is taken as an explicit hypothesis, mirroring the curved-spacetime version where no spectrum condition is available.
\end{theorem}
\begin{proof}
    \leanok
    \uses{def:cyclic-vector, thrm:von-neumann-microcausality}
    The implication is elementary ($R$ vanishes on the dense set of vectors $A\Omega$ for $A$ a local observable of $\mathbf{B}_1$, since $R(\mathbf{B}_2)$ commutes with those observables by microcausality, \ref{thrm:von-neumann-microcausality}).
\end{proof}

\begin{theorem}[Statistical Independence, bundled]
    \label{thrm:statistical-independence-bundled}
    \lean{Physicslib4.AQFT.HaagKastler.HaagKastlerNet.localVonNeumannAlgebra_separating}
    \leanfile{Physicslib4/AQFT/HaagKastler/LocalVonNeumann.lean}
    \uses{def:local-von-neumann-algebra}
    \leanok
    The separating property phrased on the bundled \texttt{VonNeumannAlgebra} $R(\mathbf{B}_2)$: with $\Omega$ cyclic for the local observables of $\mathbf{B}_1$, every $R$ in the bundled algebra $R(\mathbf{B}_2)$ of a spacelike-separated region with $R\Omega = 0$ is zero.
\end{theorem}
\begin{proof}
    \leanok
    \uses{def:local-von-neumann-algebra, thrm:statistical-independence}
    It reduces to \ref{thrm:statistical-independence} through the coercion $\uparrow R(\mathbf{B}_2) = \pi(\mathfrak{U}(\mathbf{B}_2))''$.
\end{proof}

\begin{theorem}[Additive-Free Locality via the Spacelike Complement]
    \label{thrm:additive-free-locality}
    \lean{Physicslib4.AQFT.HaagKastler.HaagKastlerNet.localVonNeumannAlgebra_le_commutant_of_subset_spacelikeComplement}
    \leanfile{Physicslib4/AQFT/HaagKastler/LocalVonNeumann.lean}
    \uses{def:spacelike-complement}
    \leanok
    Locality expressed through the spacelike complement, without attaching any algebra to the unbounded complement: for basis sets $\mathbf{B}' \subseteq \mathbf{B}^\perp$, the local von Neumann algebra $R(\mathbf{B}')$ lies in the commutant $R(\mathbf{B})'$. No local algebra is ever attached to the unbounded complement.
\end{theorem}
\begin{proof}
    \leanok
    \uses{def:spacelike-complement, thrm:von-neumann-bundled-order}
    It is a direct repackaging of bundled microcausality (\ref{thrm:von-neumann-bundled-order}) through the Galois bridge $\mathbf{B}' \subseteq \mathbf{B}^\perp \iff \mathbf{B}', \mathbf{B}$ completely spacelike, staying strictly within the physically local (bounded, diamond) regions.
\end{proof}

\begin{theorem}[Geometric Covariance of the von Neumann Net]
    \label{thrm:von-neumann-geometric-covariance}
    \lean{Physicslib4.AQFT.HaagKastler.CovariantQuasilocalAlgebra.lieConj_image_covLocalVonNeumann, Physicslib4.AQFT.HaagKastler.CovariantQuasilocalAlgebra.lieConj_image_covLocalOperators}
    \leanfile{Physicslib4/AQFT/HaagKastler/GeometricCovariance.lean}
    \uses{def:local-von-neumann, def:covariant-quasilocal-algebra}
    \leanok
    In a covariant representation $\pi$ of the quasilocal algebra with the operator covariance $U(L)\,\pi(a)\,U(L)^{-1} = \pi(\beta_L a)$, conjugation by the implementing unitary $U(L)$ carries the local von Neumann algebra of a region $\mathbf{B}$ onto that of the Lorentz-transformed region $L \cdot \mathbf{B}$:
    \[ U(L)\, R(\mathbf{B})\, U(L)^{-1} = R(L \cdot \mathbf{B}). \]
    This is the statement that the symmetry group acts \emph{geometrically} on the net of von Neumann algebras.
\end{theorem}
\begin{proof}
    \leanok
    \uses{def:local-von-neumann, def:covariant-quasilocal-algebra, thrm:irreducible-covariant-representation}
    The proof rests on three ingredients, all already available: operator covariance (the last clause of \ref{thrm:irreducible-covariant-representation}); the fact that $\beta_L = $ the covariance action sends $\iota_{\mathbf{B}}(\mathfrak{U}(\mathbf{B}))$ onto $\iota_{L\cdot\mathbf{B}}(\mathfrak{U}(L\cdot\mathbf{B}))$ (from the action-on-generators identity and surjectivity of the covariance equivalence $\alpha_L$); and the algebraic fact that conjugation by a unit is a multiplicative automorphism, which therefore commutes with the bicommutant. The reusable core is that a multiplicative automorphism maps the centralizer of a set onto the centralizer of its image, hence maps bicommutants to bicommutants. In particular $R(\mathbf{B})$ and $R(L \cdot \mathbf{B})$ are unitarily equivalent (conjugate by $U(L)$), so any unitary-conjugation-invariant property of a local von Neumann algebra --- such as being a factor --- is constant along the Lorentz orbit of a region.
\end{proof}

\begin{theorem}[Orbit-Invariance of Factoriality]
    \label{thrm:von-neumann-factor-orbit}
    \lean{Physicslib4.AQFT.HaagKastler.CovariantQuasilocalAlgebra.covLocalVonNeumann_isFactor_smul}
    \leanfile{Physicslib4/AQFT/HaagKastler/GeometricCovariance.lean}
    \uses{thrm:irreducible-factor}
    \leanok
    If the local von Neumann algebra $R(\mathbf{B})$ is a \emph{factor} (its center $R(\mathbf{B}) \cap R(\mathbf{B})'$ is exactly the scalars), then so is $R(L \cdot \mathbf{B})$ for every Lorentz transformation $L$.
\end{theorem}
\begin{proof}
    \leanok
    \uses{thrm:von-neumann-geometric-covariance}
    Geometric covariance (\ref{thrm:von-neumann-geometric-covariance}) exhibits $R(L \cdot \mathbf{B}) = U(L) R(\mathbf{B}) U(L)^{-1}$; conjugation by a unitary is a multiplicative automorphism that carries the center onto the center and fixes the scalars, so it preserves the factor property. Thus being a factor is constant along the Lorentz orbit of a region.
\end{proof}

\begin{theorem}[Geometric Covariance as a von Neumann Algebra Isomorphism]
    \label{thrm:von-neumann-covariance-iso}
    \lean{Physicslib4.AQFT.HaagKastler.CovariantQuasilocalAlgebra.covLocalVonNeumannEquiv}
    \leanfile{Physicslib4/AQFT/HaagKastler/GeometricCovariance.lean}
    \uses{def:local-von-neumann-algebra}
    \leanok
    The geometric-covariance set equality is upgraded to a first-class \emph{$*$-algebra isomorphism} of the bundled local von Neumann algebras, $R(\mathbf{B}) \cong R(L \cdot \mathbf{B})$. It is the restriction of the conjugation $*$-automorphism $T \mapsto U(L) T U(L)^{-1}$ of $\mathcal{B}(H)$ (Mathlib's \texttt{conjStarAlgEquiv}, whose adjoint-is-inverse property makes conjugation by a unitary star-preserving) to $R(\mathbf{B})$, whose image is exactly $R(L \cdot \mathbf{B})$.
\end{theorem}
\begin{proof}
    \leanok
    \uses{def:local-von-neumann-algebra, thrm:von-neumann-geometric-covariance}
    The reusable ingredient is that a $*$-automorphism carrying the underlying set of one star-subalgebra onto another restricts to a star-algebra equivalence between them. This makes the unitary equivalence of orbit-related local algebras an explicit transportable object.
\end{proof}

\subsection{Relative Commutants of Nested Local Algebras}

The theory of subalgebra inclusions $R(\mathbf{B}_1) \subseteq R(\mathbf{B}_2)$ is organised around a single object, the \emph{relative commutant} $R(\mathbf{B}_1)' \cap R(\mathbf{B}_2)$. It measures how far the smaller algebra fails to fill the larger, and always contains the center of the ambient algebra; its triviality is exactly the irreducibility of the inclusion.

\begin{definition}[Relative Commutant of a Nested Pair]
    \label{def:relative-commutant}
    \lean{Physicslib4.AQFT.HaagKastler.HaagKastlerNet.relativeCommutant}
    \uses{def:local-von-neumann-algebra, lmm:von-neumann-inter-general}
    \leanok
    For a representation $\pi$ of the quasilocal algebra and two basis regions $\mathbf{B}_1, \mathbf{B}_2$, the \emph{relative commutant} of the pair is the von Neumann algebra $R(\mathbf{B}_1)' \cap R(\mathbf{B}_2)$, the intersection of the commutant of the local von Neumann algebra $R(\mathbf{B}_1)$ with the local von Neumann algebra $R(\mathbf{B}_2)$. Since \texttt{VonNeumannAlgebra} carries no lattice meet $\sqcap$, this object is not an abstract infimum but is constructed by hand from the set $R(\mathbf{B}_1)' \cap R(\mathbf{B}_2)$, which is then its underlying set; that the intersection of two von Neumann algebras is again a von Neumann algebra is the content of \ref{lmm:von-neumann-inter-general}, applied here to the pair $R(\mathbf{B}_1)'$, $R(\mathbf{B}_2)$ (which for $\mathbf{B}_1 \neq \mathbf{B}_2$ is \emph{not} a commutant pair, so the commutant-pair form \ref{lmm:von-neumann-inter-is-von-neumann} does not suffice here). Accordingly the Lean construction rewrites the intersection as the single commutant $(R(\mathbf{B}_1)'' \cup R(\mathbf{B}_2)')'$ via \texttt{Set.centralizer\_union} and closes it under \texttt{Set.centralizer\_centralizer\_centralizer}. This is the basic object of the theory of subalgebra inclusions $R(\mathbf{B}_1) \subseteq R(\mathbf{B}_2)$.
\end{definition}

\begin{theorem}[Antitonicity of the Commutant]
    \label{thrm:commutant-antitone}
    \lean{Physicslib4.AQFT.HaagKastler.HaagKastlerNet.commutant_le_commutant_of_le}
    \leanok
    For bundled von Neumann algebras $M \le N$ on $H$, the commutants reverse the inclusion, $N' \le M'$.
\end{theorem}
\begin{proof}
    \leanok
    This is not a lattice primitive: through the coercion $\uparrow M' = \operatorname{centralizer}(\uparrow M)$ (Mathlib's \texttt{coe\_commutant}) it reduces to the antitonicity of the centralizer, $S \subseteq T \Rightarrow \operatorname{centralizer}(T) \subseteq \operatorname{centralizer}(S)$ (\texttt{Set.centralizer\_subset}), applied to $\uparrow M \subseteq \uparrow N$.
\end{proof}

\begin{theorem}[Relative Commutant Lies in the Larger Algebra]
    \label{thrm:relative-commutant-le-right}
    \lean{Physicslib4.AQFT.HaagKastler.HaagKastlerNet.relativeCommutant_le_right}
    \uses{def:relative-commutant}
    \leanok
    The relative commutant is contained in the larger algebra, $R(\mathbf{B}_1)' \cap R(\mathbf{B}_2) \le R(\mathbf{B}_2)$.
\end{theorem}
\begin{proof}
    \leanok
    \uses{def:relative-commutant}
    As \texttt{VonNeumannAlgebra} has no lattice meet, the order relation $\le$ is discharged through the coercion to underlying sets, where it is the inclusion $R(\mathbf{B}_1)' \cap R(\mathbf{B}_2) \subseteq R(\mathbf{B}_2)$, i.e. \texttt{Set.inter\_subset\_right}.
\end{proof}

\begin{theorem}[Relative Commutant Commutes with the Smaller Algebra]
    \label{thrm:relative-commutant-coe-commutant}
    \lean{Physicslib4.AQFT.HaagKastler.HaagKastlerNet.relativeCommutant_coe_subset_commutant}
    \uses{def:relative-commutant}
    \leanok
    Every element of the relative commutant commutes with all of $R(\mathbf{B}_1)$: the underlying set of $R(\mathbf{B}_1)' \cap R(\mathbf{B}_2)$ is contained in $R(\mathbf{B}_1)'$, that is $\uparrow\!\big(R(\mathbf{B}_1)' \cap R(\mathbf{B}_2)\big) \subseteq R(\mathbf{B}_1)'$.
\end{theorem}
\begin{proof}
    \leanok
    \uses{def:relative-commutant}
    On underlying sets this is \texttt{Set.inter\_subset\_left}.
\end{proof}

\begin{theorem}[Relative Commutant Contains the Center]
    \label{thrm:relative-commutant-center}
    \lean{Physicslib4.AQFT.HaagKastler.HaagKastlerNet.center_le_relativeCommutant}
    \uses{def:relative-commutant}
    \leanok
    When $\mathbf{B}_1 \subseteq \mathbf{B}_2$, the relative commutant contains the center of the ambient algebra: the underlying set of the center $R(\mathbf{B}_2) \cap R(\mathbf{B}_2)'$ is contained in that of $R(\mathbf{B}_1)' \cap R(\mathbf{B}_2)$.
\end{theorem}
\begin{proof}
    \leanok
    \uses{def:relative-commutant, thrm:von-neumann-isotony, thrm:commutant-antitone}
    Isotony (\ref{thrm:von-neumann-isotony}) gives $R(\mathbf{B}_1) \le R(\mathbf{B}_2)$; antitonicity of the commutant (\ref{thrm:commutant-antitone}) then yields $R(\mathbf{B}_2)' \le R(\mathbf{B}_1)'$, so $\uparrow R(\mathbf{B}_2)' \subseteq \uparrow R(\mathbf{B}_1)'$. Combined with $R(\mathbf{B}_2) \subseteq R(\mathbf{B}_2)$ through \texttt{Set.subset\_inter}, this lands $R(\mathbf{B}_2) \cap R(\mathbf{B}_2)'$ inside $R(\mathbf{B}_1)' \cap R(\mathbf{B}_2)$. Thus the relative commutant of a nested pair always contains the center of the ambient algebra, and is trivial (scalars) exactly when the inclusion is irreducible.
\end{proof}

\begin{definition}[Irreducible Inclusion]
    \label{def:irreducible-inclusion}
    \lean{Physicslib4.AQFT.HaagKastler.HaagKastlerNet.IsIrreducibleInclusion}
    \uses{def:relative-commutant}
    \leanok
    The inclusion $R(\mathbf{B}_1) \subseteq R(\mathbf{B}_2)$ is \emph{irreducible} when its relative commutant is trivial, $R(\mathbf{B}_1)' \cap R(\mathbf{B}_2) = \mathbb{C}\cdot 1$ (the scalar operators). This is the subfactor-theoretic notion of an irreducible inclusion; the relative commutant always contains the scalars, so irreducibility is the statement that it contains nothing more.
\end{definition}

\begin{theorem}[An Irreducible Inclusion has Factor Ambient]
    \label{thrm:irreducible-inclusion-factor}
    \lean{Physicslib4.AQFT.HaagKastler.HaagKastlerNet.isFactor_of_isIrreducibleInclusion}
    \uses{def:irreducible-inclusion}
    \leanok
    If $\mathbf{B}_1 \subseteq \mathbf{B}_2$ and the inclusion $R(\mathbf{B}_1) \subseteq R(\mathbf{B}_2)$ is irreducible, then the ambient algebra $R(\mathbf{B}_2)$ is a factor (trivial center).
\end{theorem}
\begin{proof}
    \leanok
    \uses{def:irreducible-inclusion, thrm:relative-commutant-center}
    Indeed the center $R(\mathbf{B}_2) \cap R(\mathbf{B}_2)'$ is contained in the relative commutant (\ref{thrm:relative-commutant-center}), which by irreducibility is the scalars; since the scalars are always central, the center equals the scalars. This connects irreducibility of an inclusion to factoriality of the larger algebra.
\end{proof}

\begin{theorem}[Self-Inclusion is Irreducible iff Factor]
    \label{thrm:self-inclusion-factor}
    \lean{Physicslib4.AQFT.HaagKastler.HaagKastlerNet.isIrreducibleInclusion_self_iff_isFactor}
    \uses{def:irreducible-inclusion}
    \leanok
    The trivial self-inclusion $R(\mathbf{B}) \subseteq R(\mathbf{B})$ is irreducible if and only if $R(\mathbf{B})$ is a factor.
\end{theorem}
\begin{proof}
    \leanok
    \uses{def:irreducible-inclusion}
    The relative commutant of the self-inclusion is $R(\mathbf{B})' \cap R(\mathbf{B})$, i.e. the center of $R(\mathbf{B})$ (up to the order of intersection), so its triviality is exactly the factoriality of $R(\mathbf{B})$. This is the converse-completing companion to \ref{thrm:irreducible-inclusion-factor}: irreducibility of the trivial inclusion coincides with factoriality of the algebra.
\end{proof}


\subsection{Irreducibility and Schur's Lemma}

A representation is \emph{irreducible} when its commutant is trivial: the only operators commuting with every $\pi(a)$ are scalars. This is the von Neumann (commutant) form of irreducibility. The cornerstone, connecting the commutant to the GNS state, is the topological Schur lemma for a cyclic representation.

\begin{definition}[Irreducible Representation]
    \label{def:irreducible-representation}
    \lean{Physicslib4.GNS.IsIrreducible}
    \leanfile{Physicslib4/GNS/Irreducibility.lean}
    \leanok
    A $*$-representation $\pi : A \to \mathcal{B}(H)$ is \emph{irreducible} if every bounded operator $T$ commuting with all $\pi(a)$ is a scalar multiple of the identity, $T = c \cdot 1$.
\end{definition}

\begin{theorem}[Topological Schur Lemma]
    \label{thrm:schur-lemma}
    \lean{Physicslib4.GNS.eq_smul_one_of_commute_of_cyclic}
    \leanfile{Physicslib4/GNS/Irreducibility.lean}
    \uses{def:cyclic-vector}
    \leanok
    Let $\Omega$ be a cyclic vector for $\pi$. If $T$ commutes with every $\pi(a)$ and the diagonal coefficient $a \mapsto \langle \Omega, T\,\pi(a)\Omega\rangle$ equals $c$ times $a \mapsto \langle \Omega, \pi(a)\Omega\rangle$, then $T = c \cdot 1$.
\end{theorem}
\begin{proof}
    \leanok
    \uses{def:cyclic-vector}
    The proof is pure Hilbert-space analysis: using the $*$-representation property and the commutation relation, every off-diagonal coefficient $\langle \pi(b)\Omega, (T - c)\,\pi(a)\Omega\rangle$ vanishes, so $(T - c)$ annihilates the dense cyclic orbit and is therefore zero.
\end{proof}

\begin{theorem}[Commutant Scalar iff Proportional Coefficient]
    \label{thrm:commutant-scalar-iff}
    \lean{Physicslib4.GNS.isScalar_iff_coeff_proportional}
    \leanfile{Physicslib4/GNS/Irreducibility.lean}
    \uses{def:state, thrm:gns-construction-theorem}
    \leanok
    In a cyclic representation reproducing a state $\omega$, an operator $T$ commuting with all $\pi(a)$ is a scalar multiple of the identity if and only if its diagonal coefficient $a \mapsto \langle \Omega, T\,\pi(a)\Omega\rangle$ is a scalar multiple of $\omega$. This is the precise operator-theoretic bridge to purity: irreducibility (every commutant element is scalar) is exactly the statement that every commutant coefficient is proportional to $\omega$.
\end{theorem}
\begin{proof}
    \leanok
    \uses{thrm:schur-lemma}
    This is the topological Schur lemma (\ref{thrm:schur-lemma}) read in a cyclic representation reproducing $\omega$, whose diagonal coefficient of $\pi$ is $\omega$ itself.
\end{proof}

\begin{definition}[Pure State]
    \label{def:pure-state}
    \lean{Physicslib4.GNS.IsPure}
    \leanfile{Physicslib4/GNS/Irreducibility.lean}
    \leanok
    \uses{def:state}
    A state $\omega$ is \emph{pure} if every positive linear functional $\psi$ dominated by $\omega$ (that is, $0 \le \psi(a^*a) \le \omega(a^*a)$ for all $a$) is a scalar multiple of $\omega$. This order-theoretic characterization is the extreme-point notion of purity, phrased to avoid convex-combination and normalization bookkeeping.
\end{definition}

One direction of the classical equivalence ``$\omega$ pure $\iff$ GNS representation irreducible'' is established below.

\begin{theorem}[Pure Implies Irreducible]
    \label{thrm:pure-implies-irreducible}
    \lean{Physicslib4.GNS.isIrreducible_of_isPure}
    \leanfile{Physicslib4/GNS/Irreducibility.lean}
    \uses{def:pure-state, def:irreducible-representation}
    \leanok
    If $\omega$ is pure, then any cyclic representation reproducing $\omega$ (in particular its GNS representation) is irreducible.
\end{theorem}
\begin{proof}
    \leanok
    \uses{def:pure-state, thrm:commutant-scalar-iff}
    Since the commutant is $*$-closed, a commuting operator decomposes into self-adjoint real and imaginary parts; for a self-adjoint commuting $S$, the affine rescaling $T = (2(\Vert S\Vert+1))^{-1} S + \tfrac12$ satisfies $0 \le T \le 1$, so its coefficient functional $a \mapsto \langle \Omega, T\,\pi(a)\Omega\rangle$ is positive and dominated by $\omega$. Purity forces it proportional to $\omega$, whence $T$ (and so $S$, and so the original operator) is a scalar by \ref{thrm:commutant-scalar-iff}.
\end{proof}

The converse, irreducible $\Rightarrow$ pure, requires the other half of the GNS Radon-Nikodym correspondence: that every dominated positive functional $\psi \le \omega$ is represented by an operator $T$ in the commutant, via the bounded sesquilinear form $(\pi(a)\Omega, \pi(b)\Omega) \mapsto \psi(a^* b)$. The analytic crux is that this form is well-defined and bounded.

\begin{theorem}[The GNS Radon-Nikodym Form is Bounded]
    \label{thrm:gns-form-bound}
    \lean{Physicslib4.GNS.gns_form_norm_le, Physicslib4.GNS.gns_form_well_defined}
    \leanfile{Physicslib4/GNS/RadonNikodym.lean}
    \uses{def:pure-state}
    \leanok
    For a positive functional $\psi$ dominated by the state $\omega$ in its cyclic GNS representation, the form values obey $\Vert\psi(a^* b)\Vert \le \Vert\pi(a)\Omega\Vert\,\Vert\pi(b)\Omega\Vert$. In particular the form depends only on the GNS vectors $\pi(a)\Omega, \pi(b)\Omega$, not on the representatives $a, b$, so it is well-defined on the cyclic subspace.
\end{theorem}
\begin{proof}
    \leanok
    \uses{lmm:cauchy-schwarz-inequality}
    This follows from the Cauchy-Schwarz inequality for $\psi$, the domination $\psi \le \omega$, and the reproducing identity $\omega(x^* x) = \Vert\pi(x)\Omega\Vert^2$.
\end{proof}

\begin{theorem}[The GNS Radon-Nikodym Operator]
    \label{thrm:gns-radon-nikodym-operator}
    \lean{Physicslib4.GNS.rnOp, Physicslib4.GNS.rnOp_inner, Physicslib4.GNS.rnOp_commute, Physicslib4.GNS.rnOp_reproducing}
    \leanfile{Physicslib4/GNS/RadonNikodym.lean}
    \uses{thrm:gns-form-bound}
    \leanok
    The bounded form of \ref{thrm:gns-form-bound} is the inner product of an operator $T$ on the GNS space: there is a $T$ with $\langle \pi(a)\Omega, T\,\pi(b)\Omega\rangle = \psi(a^* b)$, this $T$ commutes with every $\pi(c)$, and $\psi(a) = \langle \Omega, T\,\pi(a)\Omega\rangle$.
\end{theorem}
\begin{proof}
    \leanok
    \uses{thrm:gns-form-bound}
    Extending the densely-defined bounded form by the Fréchet-Riesz representation and a norm-controlled dense extension yields $T$ with $\langle \pi(a)\Omega, T\,\pi(b)\Omega\rangle = \psi(a^* b)$. From this reproducing identity, $T$ commutes with every $\pi(c)$ (an adjoint computation on the dense cyclic vectors), and $\psi(a) = \langle \Omega, T\,\pi(a)\Omega\rangle$ (taking $a = 1$ in the first slot).
\end{proof}

\begin{theorem}[Pure $\iff$ Irreducible]
    \label{thrm:pure-iff-irreducible}
    \lean{Physicslib4.GNS.isPure_iff_isIrreducible}
    \leanfile{Physicslib4/GNS/RadonNikodym.lean}
    \leanok
    A state $\omega$ is pure if and only if its cyclic GNS representation is irreducible.
\end{theorem}
\begin{proof}
    \leanok
    \uses{thrm:pure-implies-irreducible, thrm:gns-radon-nikodym-operator, thrm:commutant-scalar-iff}
    For the converse direction, given a dominated $\psi$, its Radon-Nikodym operator $T$ commutes with $\pi$, so by irreducibility $T = c \cdot 1$; the reproducing identity then gives $\psi(a) = c\,\omega(a)$, so $\psi$ is a scalar multiple of $\omega$ and $\omega$ is pure.
\end{proof}

Irreducibility has a von Neumann algebraic reading: the representation generates a \emph{factor}.

\begin{theorem}[An Irreducible Representation Generates a Factor]
    \label{thrm:irreducible-factor}
    \lean{Physicslib4.GNS.center_gnsVonNeumann_eq_of_isIrreducible}
    \leanfile{Physicslib4/GNS/Irreducibility.lean}
    \uses{def:irreducible-representation, def:local-von-neumann}
    \leanok
    The von Neumann algebra $\pi(A)''$ generated by an irreducible representation has \emph{trivial center}: an operator lies in the center $\pi(A)'' \cap (\pi(A)'')'$ if and only if it is a scalar multiple of the identity.
\end{theorem}
\begin{proof}
    \leanok
    \uses{def:irreducible-representation, def:local-von-neumann}
    Indeed the center is contained in $(\pi(A)'')' = \pi(A)'$ (the triple commutant collapses to the single one), which irreducibility makes the scalars; conversely scalars are central.
\end{proof}

\begin{theorem}[Irreducibility $\iff$ Generating $\mathcal{B}(H)$]
    \label{thrm:irreducible-generates-all}
    \lean{Physicslib4.GNS.isIrreducible_iff_gnsVonNeumann_eq_univ, Physicslib4.GNS.gnsVonNeumann_eq_univ_of_isIrreducible}
    \leanfile{Physicslib4/GNS/Irreducibility.lean}
    \uses{def:irreducible-representation, def:local-von-neumann}
    \leanok
    A representation is irreducible if and only if the von Neumann algebra $\pi(A)''$ it generates is \emph{all} of $\mathcal{B}(H)$: $\pi(A)'' = \mathcal{B}(H)$. This is the density (bicommutant-theorem) form of irreducibility, sharpening the factor statement \ref{thrm:irreducible-factor}.
\end{theorem}
\begin{proof}
    \leanok
    \uses{def:irreducible-representation, def:local-von-neumann}
    For the forward direction, irreducibility makes the commutant $\pi(A)'$ exactly the scalars, and the centralizer of the scalars is everything (every bounded operator commutes with $c \cdot 1$), so $\pi(A)'' = (\text{scalars})' = \mathcal{B}(H)$. The converse uses that $\mathcal{B}(H)$ is a \emph{central} $\mathbb{C}$-algebra (its center is the scalars): if the bicommutant is everything, every operator commutes with the commutant, so any operator commuting with $\pi(A)$ is central, hence scalar.
\end{proof}

\begin{theorem}[Bundled Density Form of Irreducibility]
    \label{thrm:gns-generates-all-bundled}
    \lean{Physicslib4.GNS.coe_gnsVonNeumannAlgebra_eq_univ_of_isIrreducible, Physicslib4.GNS.gnsVonNeumannAlgebra_isGreatest_of_isIrreducible}
    \leanfile{Physicslib4/GNS/Irreducibility.lean}
    \uses{def:local-von-neumann-algebra, lmm:bicommutant-of-selfadjoint-is-von-neumann}
    \leanok
    The von Neumann algebra $\pi(A)''$ generated by a representation is bundled as a first-class \texttt{VonNeumannAlgebra} (\texttt{gnsVonNeumannAlgebra}); the bundling is \ref{lmm:bicommutant-of-selfadjoint-is-von-neumann} applied to the self-adjoint image $\pi(A)$. For an irreducible representation, this bundled algebra is all of $\mathcal{B}(H)$: its underlying set is everything, and equivalently it is the \emph{greatest} von Neumann algebra on $H$ (every $S$ satisfies $S \le \pi(A)''$). Mathlib's \texttt{VonNeumannAlgebra} carries no lattice $\top$, so the literal $\pi(A)'' = \top$ is phrased at the level of the underlying set together with the greatest-element statement under the existing $\le$.
\end{theorem}
\begin{proof}
    \leanok
    \uses{thrm:irreducible-generates-all, lmm:bicommutant-of-selfadjoint-is-von-neumann}
    The bundling is \ref{lmm:bicommutant-of-selfadjoint-is-von-neumann} for the self-adjoint set $\pi(A)$, and the identification of the underlying set is the bundled counterpart of \ref{thrm:irreducible-generates-all}.
\end{proof}

\begin{theorem}[The GNS Representation of a Pure State is a Factor]
    \label{thrm:pure-factor}
    \lean{Physicslib4.GNS.exists_gns_factor_of_isPure}
    \leanfile{Physicslib4/GNS/RadonNikodym.lean}
    \uses{def:pure-state, thrm:gns-construction-theorem}
    \leanok
    For a pure state $\omega$ there is a cyclic GNS triple $(H, \pi, \Omega)$ reproducing $\omega$ whose generated von Neumann algebra $\pi(A)''$ has trivial center, i.e. is a factor.
\end{theorem}
\begin{proof}
    \leanok
    \uses{thrm:pure-iff-irreducible, thrm:irreducible-factor, thrm:gns-construction-theorem}
    This combines the GNS construction, purity $\Rightarrow$ irreducibility (\ref{thrm:pure-iff-irreducible}), and \ref{thrm:irreducible-factor}. It applies verbatim to the quasilocal algebra $\mathfrak{U}$ and to each local algebra $\mathfrak{U}(\mathbf{B})$, since both are C*-algebras.
\end{proof}

\begin{theorem}[The GNS Representation of a Pure State Generates $\mathcal{B}(H)$]
    \label{thrm:pure-generates-all}
    \lean{Physicslib4.GNS.exists_gns_generates_all_of_isPure}
    \leanfile{Physicslib4/GNS/RadonNikodym.lean}
    \uses{def:pure-state, thrm:gns-construction-theorem}
    \leanok
    The density (bicommutant-theorem) sharpening of \ref{thrm:pure-factor}: for a pure state $\omega$ there is a cyclic GNS triple $(H, \pi, \Omega)$ reproducing $\omega$ whose generated von Neumann algebra is \emph{all} of $\mathcal{B}(H)$, $\pi(A)'' = \mathcal{B}(H)$.
\end{theorem}
\begin{proof}
    \leanok
    \uses{thrm:pure-iff-irreducible, thrm:irreducible-generates-all, thrm:gns-construction-theorem}
    It combines the GNS construction, purity $\Rightarrow$ irreducibility (\ref{thrm:pure-iff-irreducible}), and the density form \ref{thrm:irreducible-generates-all}. As with the factor statement, it applies verbatim to the quasilocal algebra $\mathfrak{U}$ and to each local algebra $\mathfrak{U}(\mathbf{B})$.
\end{proof}

The order-theoretic definition of purity (\ref{def:pure-state}) coincides with the convex-geometric one: $\omega$ is an extreme point of the state space. The analytic key is the normalization identity for positive functionals.

\begin{theorem}[Norm of a Positive Functional]
    \label{thrm:norm-positive-functional}
    \lean{Physicslib4.GNS.norm_eq_re_apply_one_of_positive}
    \leanfile{Physicslib4/GNS/ExtremeState.lean}
    \leanok
    For a positive linear functional $\varphi$ on a unital C*-algebra, $\Vert\varphi\Vert = \mathrm{Re}\,\varphi(1)$. In particular every state satisfies $\omega(1) = 1$.
\end{theorem}
\begin{proof}
    \leanok
    \uses{lmm:cauchy-schwarz-inequality}
    The bound $\mathrm{Re}\,\varphi(1) \le \Vert\varphi\Vert$ is immediate from $\Vert 1\Vert = 1$; conversely Cauchy-Schwarz with the first slot equal to $1$ gives $\Vert\varphi(b)\Vert^2 \le \mathrm{Re}\,\varphi(1)\cdot\mathrm{Re}\,\varphi(b^* b) \le \mathrm{Re}\,\varphi(1)\cdot\Vert\varphi\Vert\,\Vert b\Vert^2$, whence $\Vert\varphi\Vert^2 \le \mathrm{Re}\,\varphi(1)\cdot\Vert\varphi\Vert$.
\end{proof}

\begin{definition}[Extreme Point of the State Space]
    \label{def:extreme-state}
    \lean{Physicslib4.GNS.State.IsExtremePoint}
    \leanfile{Physicslib4/GNS/ExtremeState.lean}
    \leanok
    \uses{def:state}
    A state $\omega$ is an \emph{extreme point} of the (convex) state space if it is not a nontrivial convex combination of two distinct states: whenever $\omega = t\,\omega_1 + (1-t)\,\omega_2$ with $0 < t < 1$ and $\omega_1, \omega_2$ states, then $\omega_1 = \omega_2$.
\end{definition}

\begin{theorem}[Pure $\iff$ Extreme Point]
    \label{thrm:pure-iff-extreme}
    \lean{Physicslib4.GNS.isPure_iff_isExtremePoint}
    \leanfile{Physicslib4/GNS/ExtremeState.lean}
    \uses{def:pure-state, def:extreme-state}
    \leanok
    A state $\omega$ is pure if and only if it is an extreme point of the state space.
\end{theorem}
\begin{proof}
    \leanok
    \uses{def:pure-state, def:extreme-state, thrm:norm-positive-functional}
    If $\omega$ is pure and $\omega = t\,\omega_1 + (1-t)\,\omega_2$, then $t\,\omega_1$ is a positive functional dominated by $\omega$, hence (by purity) a scalar multiple of $\omega$; evaluating at $1$, where every state gives $1$, pins the scalar to $t$ and forces $\omega_1 = \omega = \omega_2$. Conversely, if $\omega$ is extreme and $\psi \le \omega$ is dominated, set $\lambda = \mathrm{Re}\,\psi(1) \in [0,1]$; for $\lambda \in (0,1)$ the rescaled functionals $\lambda^{-1}\psi$ and $(1-\lambda)^{-1}(\omega - \psi)$ are states (normalized via \ref{thrm:norm-positive-functional}) whose convex combination is $\omega$, so extremality identifies them and forces $\psi = \lambda\,\omega$; the boundary cases $\lambda = 0$ and $\lambda = 1$ give $\psi = 0$ and $\psi = \omega$ respectively.
\end{proof}

\begin{theorem}[State-Space Convexity and the Extreme-Point Bridge]
    \label{thrm:state-space-convex-bridge}
    \lean{Physicslib4.GNS.convex_stateSpace, Physicslib4.GNS.mem_extremePoints_iff_isExtremePoint}
    \leanfile{Physicslib4/GNS/ExtremeState.lean}
    \uses{def:state, def:extreme-state}
    \leanok
    Realizing the state space as the subset $\mathrm{stateSpace}(A) \subseteq A \to_L \mathbb{C}$ (the latter a real topological vector space via $\texttt{NormedSpace.complexToReal}$), it is convex: a real convex combination $a\,\omega_1 + b\,\omega_2$ of states is again a state, via $\texttt{State.convexCombo}$. Moreover a state $\omega$ lies in Mathlib's $\mathrm{extremePoints}_{\mathbb{R}}$ of the state space exactly when it is extreme in the sense of \ref{def:extreme-state}, connecting the purity characterizations to the Krein-Milman/Choquet API.
\end{theorem}
\begin{proof}
    \leanok
    \uses{def:state, def:extreme-state}
    Convexity is witnessed by \texttt{State.convexCombo}, which exhibits a real convex combination of two states as a state; the extreme-point bridge is the unfolding of Mathlib's $\mathrm{extremePoints}_{\mathbb{R}}$ against \ref{def:extreme-state}.
\end{proof}

\begin{definition}[Pullback of a State]
    \label{def:state-pullback}
    \lean{Physicslib4.GNS.State.comp}
    \uses{def:state}
    \leanok
    For a unital $*$-homomorphism $\pi : A \to B$ of C*-algebras and a state $\omega$ on $B$, the \emph{pullback} $\omega \circ \pi$ is the functional $a \mapsto \omega(\pi(a))$. It is again a state: positivity is the $*$-compatibility $\omega(\pi(a^* a)) = \omega(\pi(a)^*\pi(a)) \ge 0$, and normalization $\Vert\omega\circ\pi\Vert = 1$ follows from unitality $\pi(1) = 1$ together with \ref{thrm:norm-positive-functional}. Thus $A \mapsto \mathrm{State}(A)$ is contravariant in $A$: a $*$-homomorphism $\pi : A \to B$ induces the pullback $\mathrm{State}(B) \to \mathrm{State}(A)$.
\end{definition}

\begin{theorem}[Functoriality of the State Pullback]
    \label{thrm:state-pullback-functorial}
    \lean{Physicslib4.GNS.State.comp_id, Physicslib4.GNS.State.comp_comp}
    \uses{def:state-pullback}
    \leanok
    The pullback is a contravariant functor on C*-algebras: pulling back along the identity $*$-homomorphism is the identity ($\omega \circ \mathrm{id} = \omega$), and pulling back along a composite reverses order, $\omega \circ (\pi_2 \circ \pi_1) = (\omega \circ \pi_2) \circ \pi_1$ for $\pi_1 : A \to B$, $\pi_2 : B \to C$ and $\omega$ a state on $C$.
\end{theorem}
\begin{proof}
    \leanok
    \uses{def:state-pullback}
    Both identities are immediate from the defining equation $(\omega \circ \pi)(a) = \omega(\pi(a))$ of the pullback (\ref{def:state-pullback}), checked pointwise on $a$.
\end{proof}

\begin{theorem}[Purity is Invariant under a $*$-Isomorphism]
    \label{thrm:pure-pullback-invariant}
    \lean{Physicslib4.GNS.isPure_comp_iff}
    \uses{def:state-pullback, def:pure-state}
    \leanok
    For a $*$-isomorphism $\Phi : A \simeq B$ of C*-algebras and a state $\omega$ on $B$, the pullback $\omega \circ \Phi$ is pure if and only if $\omega$ is. This is the cross-algebra generalization of the $*$-automorphism case (purity is a covariance invariant); a $*$-isomorphism identifies the pure states of $A$ and $B$.
\end{theorem}
\begin{proof}
    \leanok
    \uses{def:state-pullback, def:pure-state}
    A dominated positive functional $\psi \le \omega \circ \Phi$ transports to $\psi \circ \Phi^{-1} \le \omega$, which purity sends to a scalar multiple of $\omega$; transporting back gives $\psi$ proportional to $\omega \circ \Phi$, and the converse follows by applying the same to $\Phi^{-1}$ (using $(\omega \circ \Phi)\circ \Phi^{-1} = \omega$).
\end{proof}

\begin{theorem}[Weak-* Compactness of the State Space]
    \label{thrm:state-space-weak-compact}
    \lean{Physicslib4.GNS.isCompact_weakStateSet}
    \leanfile{Physicslib4/GNS/PureStateExists.lean}
    \uses{def:state}
    \leanok
    The state space, realized inside the weak-* dual $\mathrm{WeakDual}\,\mathbb{C}\,A$ as the positive functionals with $\varphi(1) = 1$, is \emph{weak-* compact}. This is the analytic input for the existence of pure states via Krein-Milman.
\end{theorem}
\begin{proof}
    \leanok
    \uses{def:state}
    By Banach-Alaoglu it suffices that it is a weak-* closed subset of the closed unit ball: the positivity conditions $0 \le \varphi(a^* a)$ and the normalization $\varphi(1) = 1$ are each weak-* closed (evaluation at a fixed element is weak-* continuous), and a positive functional with $\varphi(1)=1$ has $\|\varphi\| = 1$.
\end{proof}

These equivalences specialize to the quasilocal algebra $\mathfrak{U}$ of a Minkowski net, where they characterize purity of a global state - the natural setting for the vacuum and other distinguished states.

\begin{theorem}[Pure $\iff$ Extreme Point on the Quasilocal Algebra]
    \label{thrm:pure-iff-extreme-quasilocal}
    \lean{Physicslib4.AQFT.HaagKastler.HaagKastlerNet.pure_iff_extreme}
    \leanfile{Physicslib4/AQFT/HaagKastler/Purity.lean}
    \uses{def:pure-state, def:extreme-state, def:quasilocal-completeness}
    \leanok
    A state $\omega$ on the canonical quasilocal algebra $\mathfrak{U}$ of a Haag-Kastler net is pure if and only if it is an extreme point of the state space of $\mathfrak{U}$.
\end{theorem}
\begin{proof}
    \leanok
    \uses{thrm:pure-iff-extreme}
    This is the abstract equivalence \ref{thrm:pure-iff-extreme} applied to the C*-algebra $\mathfrak{U}$.
\end{proof}

\begin{theorem}[Pure $\iff$ Irreducible GNS on the Quasilocal Algebra]
    \label{thrm:pure-iff-irreducible-quasilocal}
    \lean{Physicslib4.AQFT.HaagKastler.HaagKastlerNet.exists_gns_pure_iff_irreducible}
    \leanfile{Physicslib4/AQFT/HaagKastler/Purity.lean}
    \uses{def:pure-state, def:irreducible-representation, thrm:gns-construction-theorem, def:quasilocal-completeness}
    \leanok
    For a state $\omega$ on the quasilocal algebra $\mathfrak{U}$ there is a GNS triple $(H, \pi, \Omega)$ reproducing $\omega$ in which $\omega$ is pure if and only if the representation $\pi$ is irreducible.
\end{theorem}
\begin{proof}
    \leanok
    \uses{thrm:pure-iff-irreducible, thrm:gns-construction-theorem}
    This combines the GNS construction with the abstract \ref{thrm:pure-iff-irreducible}.
\end{proof}

\subsection{Unitary Equivalence and Superselection}

The comparison of representations is the entry point to superselection theory. For $*$-representations $\pi : A \to \mathcal{B}(H)$ of a C*-algebra --- in particular of the quasilocal algebra $\mathfrak{U}$ or a local algebra $\mathfrak{U}(\mathbf{B})$ --- we record three successively coarser notions of sameness (unitary equivalence, quasi-equivalence, and the negation of the finest, disjointness) and the representation-theoretic invariants they preserve.

\begin{definition}[Unitary Equivalence of Representations]
    \label{def:unitary-equivalence}
    \lean{Physicslib4.GNS.UnitaryEquiv}
    \leanfile{Physicslib4/GNS/UnitaryEquiv.lean}
    \leanok
    \uses{thrm:gns-construction-theorem}
    Two $*$-representations $\pi_1 : A \to \mathcal{B}(H_1)$ and $\pi_2 : A \to \mathcal{B}(H_2)$ are \emph{unitarily equivalent} when there is an isometric isomorphism $U : H_1 \simeq H_2$ of the underlying Hilbert spaces intertwining them: $U(\pi_1(a)\,x) = \pi_2(a)(U x)$. This is an equivalence relation --- reflexive, symmetric, and transitive.
\end{definition}

\begin{theorem}[Irreducibility and Factoriality are Unitary Invariants]
    \label{thrm:unitary-equiv-invariants}
    \lean{Physicslib4.GNS.UnitaryEquiv.isIrreducible_iff, Physicslib4.GNS.UnitaryEquiv.isFactor_iff}
    \leanfile{Physicslib4/GNS/UnitaryEquiv.lean}
    \uses{def:unitary-equivalence, def:irreducible-representation, thrm:irreducible-factor}
    \leanok
    Unitarily equivalent representations share their representation-theoretic type: $\pi_1$ is irreducible if and only if $\pi_2$ is, and the generated von Neumann algebra $\pi_1(A)''$ is a factor if and only if $\pi_2(A)''$ is.
\end{theorem}
\begin{proof}
    \leanok
    \uses{def:unitary-equivalence, def:irreducible-representation, thrm:irreducible-factor}
    The transport is packaged through the cross-space conjugation $T \mapsto U T U^{-1}$, a multiplicative isomorphism $\mathcal{B}(H_1) \simeq \mathcal{B}(H_2)$ that carries $\pi_1(A)$ onto $\pi_2(A)$, centralizers onto centralizers, and scalars onto scalars; irreducibility (commutant equal to the scalars) and factoriality (trivial center of $\pi(A)''$) are then preserved.
\end{proof}

\subsection{GNS Covariance}

\begin{lemma}[Cyclicity pulls back along a surjective $*$-homomorphism]
    \label{lmm:cyclic-pullback-surjective}
    \lean{Physicslib4.GNS.isCyclicVector_comp_of_surjective}
    \uses{def:cyclic-vector, def:state-pullback}
    \leanok
    Let $\pi : B \to \mathcal{B}(H)$ be a $*$-representation with cyclic vector $\Omega$, and let $\Phi : A \to B$ be a surjective unital $*$-homomorphism. Here the pulled-back representation $\pi \circ \Phi$ of $A$ is the composite $*$-homomorphism (Mathlib's \texttt{StarAlgHom.comp}), just as the pullback of a state is fixed by \ref{def:state-pullback}. Then $\Omega$ is cyclic for $\pi \circ \Phi$.
\end{lemma}
\begin{proof}
    \leanok
    \uses{def:cyclic-vector}
    Indeed surjectivity of $\Phi$ gives $\{\pi(\Phi(a))\Omega : a \in A\} = \{\pi(b)\Omega : b \in B\}$, so the two orbits coincide as sets and density transfers verbatim.
\end{proof}

\begin{theorem}[GNS covariance under a $*$-isomorphism]
    \label{thrm:gns-covariance}
    \lean{Physicslib4.GNS.exists_unitary_of_gns_comp}
    \uses{def:state-pullback, def:cyclic-vector, thrm:gns-construction-theorem}
    \leanok
    Let $\Phi : A \simeq B$ be a $*$-isomorphism of unital C*-algebras and $\omega$ a state on $B$. Let $(H_1, \pi_1, \Omega_1)$ be a cyclic representation of $A$ reproducing the pullback state $\omega \circ \Phi$, and $(H_2, \pi_2, \Omega_2)$ a cyclic representation of $B$ reproducing $\omega$. Then there is a unitary $U : H_1 \simeq H_2$ with $U\Omega_1 = \Omega_2$ intertwining the two representations along $\Phi$: $U(\pi_1(a)x) = \pi_2(\Phi(a))(Ux)$ for all $a \in A$, $x \in H_1$. In other words the GNS data transports covariantly along an isomorphism of the algebras.
\end{theorem}
\begin{proof}
    \leanok
    \uses{lmm:cyclic-pullback-surjective, thrm:gns-construction-theorem, def:state-pullback}
    The proof observes that $(H_2, \pi_2 \circ \Phi, \Omega_2)$ is itself a cyclic representation of $A$ (\ref{lmm:cyclic-pullback-surjective}, $\Phi$ being surjective) and that it reproduces $\omega \circ \Phi$, since $\langle \Omega_2, \pi_2(\Phi a)\Omega_2\rangle = \omega(\Phi a) = (\omega \circ \Phi)(a)$ by the defining equation of the pullback state. One then applies the \emph{uniqueness clause} of \ref{thrm:gns-construction-theorem} --- in Lean the separate declaration \texttt{Physicslib4.GNS.gns\_unique}, not \texttt{gns\_construction} --- to the two cyclic representations of $A$ attached to the single state $\omega \circ \Phi$.
\end{proof}

\begin{theorem}[The GNS representation of a pullback state]
    \label{thrm:gns-covariance-unitary-equiv}
    \lean{Physicslib4.GNS.unitaryEquiv_comp_of_gns}
    \uses{def:unitary-equivalence, def:state-pullback}
    \leanok
    Restated in the language of unitary equivalence. Let $\Phi : A \simeq B$ be a $*$-isomorphism of unital C*-algebras and $\omega$ a state on $B$; let $(H_1, \pi_1, \Omega_1)$ be a cyclic representation of $A$ reproducing $\omega \circ \Phi$ and $(H_2, \pi_2, \Omega_2)$ a cyclic representation of $B$ reproducing $\omega$. Then $\pi_1$ and $\pi_2 \circ \Phi$ are unitarily equivalent (\ref{def:unitary-equivalence}).
\end{theorem}
\begin{proof}
    \leanok
    \uses{thrm:gns-covariance, def:unitary-equivalence}
    The intertwining unitary of \ref{thrm:gns-covariance} is exactly the required witness. Together with the invariance of irreducibility and factoriality (\ref{thrm:unitary-equiv-invariants}), this is the mechanism by which superselection sectors transport along an isomorphism of the observable algebra.
\end{proof}

\begin{lemma}[Pullback along a surjection preserves the image algebra]
    \label{lmm:pullback-image-invariants}
    \lean{Physicslib4.GNS.range_comp_of_surjective, Physicslib4.GNS.isIrreducible_comp_iff, Physicslib4.GNS.gnsVonNeumann_comp_of_surjective}
    \uses{def:irreducible-representation, def:local-von-neumann}
    \leanok
    Let $\pi : B \to \mathcal{B}(H)$ be a $*$-representation and $\Phi : A \to B$ a surjective unital $*$-homomorphism, and let $\pi \circ \Phi$ be the composite $*$-homomorphism (Mathlib's \texttt{StarAlgHom.comp}). Then $\pi \circ \Phi$ has the same image as $\pi$: $(\pi \circ \Phi)(A) = \pi(B)$. Consequently everything computed from the image alone is unchanged: $\pi \circ \Phi$ is irreducible if and only if $\pi$ is (irreducibility being triviality of the commutant of the image), and the generated von Neumann algebras (\ref{def:local-von-neumann}) coincide, $(\pi \circ \Phi)(A)'' = \pi(B)''$. The conclusion is sharper than a mere unitary equivalence would give: no transport is involved at all, since the two representations act on the \emph{same} Hilbert space and generate literally the same algebras, not just isomorphic ones.
\end{lemma}
\begin{proof}
    \leanok
    \uses{def:irreducible-representation, def:local-von-neumann}
    The image identity is surjectivity of $\Phi$ applied to ranges, $\operatorname{range}(\pi \circ \Phi) = \operatorname{range}(\pi)$, which is Mathlib's \texttt{Function.Surjective.range\_comp}. For irreducibility, the existing repository lemma \texttt{Physicslib4.GNS.isIrreducible\_iff\_centralizer} (in \texttt{Physicslib4/GNS/UnitaryEquiv.lean}) already identifies irreducibility with triviality of the centralizer of the image, $\operatorname{centralizer}(\operatorname{range} \pi) = \mathbb{C} \cdot 1$; rewriting that criterion along the image identity gives the equivalence with no re-derivation. The generated von Neumann algebra is the double commutant of the image, so the same rewriting gives $(\pi \circ \Phi)(A)'' = \pi(B)''$.
\end{proof}

\begin{theorem}[Superselection type transports along a $*$-isomorphism]
    \label{thrm:gns-sector-transport}
    \lean{Physicslib4.GNS.isIrreducible_iff_of_gns_comp, Physicslib4.GNS.isFactor_iff_of_gns_comp}
    \uses{thrm:gns-covariance-unitary-equiv, thrm:unitary-equiv-invariants, lmm:pullback-image-invariants, def:state-pullback}
    \leanok
    Let $\Phi : A \simeq B$ be a $*$-isomorphism of unital C*-algebras and $\omega$ a state on $B$; let $(H_1, \pi_1, \Omega_1)$ be a cyclic representation of $A$ reproducing $\omega \circ \Phi$ and $(H_2, \pi_2, \Omega_2)$ a cyclic representation of $B$ reproducing $\omega$. Then $\pi_1$ is irreducible if and only if $\pi_2$ is, and $\pi_1(A)''$ is a factor if and only if $\pi_2(B)''$ is. The chain is: $\pi_1$ is unitarily equivalent to $\pi_2 \circ \Phi$ (\ref{thrm:gns-covariance-unitary-equiv}), unitary equivalence preserves irreducibility and factoriality (\ref{thrm:unitary-equiv-invariants}), and $\pi_2 \circ \Phi$ has the same image algebra as $\pi_2$ (\ref{lmm:pullback-image-invariants}). Thus an isomorphism of the observable algebra carries superselection sectors to superselection sectors, preserving their type; together with the invariance of purity (\ref{thrm:pure-pullback-invariant}) it shows the whole sector structure is an invariant of the algebra, not of its presentation.
\end{theorem}
\begin{proof}
    \leanok
    \uses{thrm:gns-covariance-unitary-equiv, thrm:unitary-equiv-invariants, lmm:pullback-image-invariants}
    Compose the three steps. \ref{thrm:gns-covariance-unitary-equiv} gives $\pi_1 \sim_u \pi_2 \circ \Phi$; \ref{thrm:unitary-equiv-invariants} turns this into the two equivalences $\mathrm{Irr}(\pi_1) \iff \mathrm{Irr}(\pi_2 \circ \Phi)$ and $\mathrm{Factor}(\pi_1(A)'') \iff \mathrm{Factor}((\pi_2 \circ \Phi)(A)'')$; and \ref{lmm:pullback-image-invariants}, applied to the surjection $\Phi$, replaces $\pi_2 \circ \Phi$ by $\pi_2$ on the right of each.
\end{proof}

\begin{theorem}[GNS covariance for local algebras]
    \label{thrm:gns-covariance-local}
    \lean{Physicslib4.AQFT.HaagKastler.HaagKastlerNet.unitaryEquiv_gns_covEquiv, Physicslib4.AQFT.HaagKastler.HaagKastlerNet.isIrreducible_iff_gns_covEquiv, Physicslib4.AQFT.HaagKastler.HaagKastlerNet.isFactor_iff_gns_covEquiv}
    \uses{def:haag-kastler-net, def:lorentz-covariance, thrm:gns-covariance-unitary-equiv, thrm:gns-sector-transport, def:state-pullback}
    \leanok
    Let $\mathfrak{U}$ be a Haag-Kastler net (\ref{def:haag-kastler-net}). The covariance equivalence $\alpha_L : \mathfrak{U}(\mathbf{B}) \simeq \mathfrak{U}(L\cdot\mathbf{B})$ supplied by Axiom 5 (\ref{def:lorentz-covariance}) is a $*$-isomorphism of local algebras, so the abstract GNS-covariance results apply to it region by region. Let $\omega$ be a state on $\mathfrak{U}(L\cdot\mathbf{B})$, so that $\omega \circ \alpha_L$ is a state on $\mathfrak{U}(\mathbf{B})$ (\ref{def:state-pullback}); let $(H_1, \pi_1, \Omega_1)$ be a cyclic representation of $\mathfrak{U}(\mathbf{B})$ reproducing $\omega \circ \alpha_L$ and $(H_2, \pi_2, \Omega_2)$ a cyclic representation of $\mathfrak{U}(L\cdot\mathbf{B})$ reproducing $\omega$. Then $\pi_1$ is unitarily equivalent to $\pi_2 \circ \alpha_L$; moreover $\pi_1$ is irreducible if and only if $\pi_2$ is, and $\pi_1(\mathfrak{U}(\mathbf{B}))''$ is a factor if and only if $\pi_2(\mathfrak{U}(L\cdot\mathbf{B}))''$ is. In physical terms: the superselection type of a local state is constant along the Lorentz orbit of the region.
\end{theorem}
\begin{proof}
    \leanok
    \uses{thrm:gns-covariance-unitary-equiv, thrm:gns-sector-transport}
    Instantiate \ref{thrm:gns-covariance-unitary-equiv} and \ref{thrm:gns-sector-transport} at $A := \mathfrak{U}(\mathbf{B})$, $B := \mathfrak{U}(L\cdot\mathbf{B})$ and $\Phi := \alpha_L$, the $*$-isomorphism extracted from the net's Lorentz covariance datum.
\end{proof}

\subsection{Disjointness and Quasi-Equivalence}

\begin{definition}[Disjoint Representations]
    \label{def:disjoint-representations}
    \lean{Physicslib4.GNS.AreDisjoint}
    \leanfile{Physicslib4/GNS/Superselection.lean}
    \leanok
    \uses{def:unitary-equivalence}
    An \emph{intertwiner} $T : H_1 \to H_2$ between $\pi_1$ and $\pi_2$ satisfies $T(\pi_1(a)\,x) = \pi_2(a)(T x)$; intertwiners are closed under sums, scalars, composition, and the adjoint (the adjoint of a $\pi_1 \to \pi_2$ intertwiner is a $\pi_2 \to \pi_1$ intertwiner). Two representations are \emph{disjoint} when the only intertwiner between them is $0$. Disjointness is symmetric (take adjoints), and a representation on a nonzero Hilbert space is never disjoint from itself; more generally, unitarily equivalent representations on nonzero spaces are never disjoint, since the implementing unitary is a nonzero intertwiner.
\end{definition}

\begin{definition}[Quasi-Equivalence of Representations]
    \label{def:quasi-equivalence}
    \lean{Physicslib4.GNS.QuasiEquiv}
    \leanfile{Physicslib4/GNS/Superselection.lean}
    \leanok
    \uses{def:unitary-equivalence, thrm:irreducible-factor}
    Two representations are \emph{quasi-equivalent} when there is a $*$-isomorphism of their generated von Neumann algebras $\pi_1(A)'' \simeq \pi_2(A)''$ carrying $\pi_1(a)$ to $\pi_2(a)$. This is an equivalence relation, and it is coarser than unitary equivalence: unitary equivalence implies quasi-equivalence, because the conjugation $*$-isomorphism $T \mapsto U T U^{-1}$ restricts to a $*$-isomorphism of the generated von Neumann algebras.
\end{definition}

\begin{theorem}[Schur's Lemma and the Irreducible Dichotomy]
    \label{thrm:irreducible-dichotomy}
    \lean{Physicslib4.GNS.UnitaryEquiv.of_intertwines_of_isIrreducible, Physicslib4.GNS.areDisjoint_or_unitaryEquiv_of_isIrreducible}
    \leanfile{Physicslib4/GNS/Superselection.lean}
    \uses{def:unitary-equivalence, def:disjoint-representations, def:irreducible-representation}
    \leanok
    A nonzero intertwiner $T$ between two \emph{irreducible} representations rescales to a unitary equivalence. Consequently two irreducible representations are either \emph{disjoint} or \emph{unitarily equivalent} --- the foundational trichotomy of superselection theory, identifying sectors with unitary-equivalence classes of irreducible representations.
\end{theorem}
\begin{proof}
    \leanok
    \uses{def:irreducible-representation, def:unitary-equivalence, def:disjoint-representations}
    $T^\ast T$ commutes with $\pi_1$, hence is a positive scalar $r \cdot 1$ with $r > 0$, so $(\sqrt r)^{-1}\,T$ is a linear isometry; and $T T^\ast$ is a nonzero scalar (irreducibility of $\pi_2$), which makes $T$ surjective, so the isometry is a unitary.
\end{proof}

\begin{lemma}[Schur Multiplicity]
    \label{lmm:schur-multiplicity}
    \lean{Physicslib4.GNS.eq_smul_of_intertwines_of_isIrreducible}
    \leanfile{Physicslib4/GNS/Superselection.lean}
    \uses{def:irreducible-representation}
    \leanok
    The space of intertwiners between two irreducible representations is at most one-dimensional: any two intertwiners $S, T$ with $S \neq 0$ are proportional, $T = \lambda \cdot S$.
\end{lemma}
\begin{proof}
    \leanok
    \uses{thrm:irreducible-dichotomy, def:irreducible-representation}
    Indeed $S^\ast S = a\cdot 1$ and $S^\ast T = b \cdot 1$ (commuting with $\pi_1$, hence scalar), and $S S^\ast = c \cdot 1$ with $c \neq 0$ (irreducibility of $\pi_2$) makes $S^\ast$ injective; since $S^\ast\big(T - (b/a)\,S\big) = 0$, injectivity gives $T = (b/a)\,S$. So irreducible representations intertwine multiplicity-free.
\end{proof}

\begin{lemma}[Endomorphism Algebra of an Irreducible Representation]
    \label{lmm:endomorphism-scalar}
    \lean{Physicslib4.GNS.intertwines_self_iff_isScalar}
    \leanfile{Physicslib4/GNS/Superselection.lean}
    \uses{def:irreducible-representation, def:disjoint-representations}
    \leanok
    The endomorphism algebra of an irreducible representation is $\mathbb{C}\cdot 1$: every self-intertwiner of an irreducible representation is a scalar multiple of the identity.
\end{lemma}
\begin{proof}
    \leanok
    \uses{def:irreducible-representation, def:disjoint-representations}
    This is the commutant form of irreducibility read through the intertwiner language, and it identifies each irreducible sector with a simple (one-dimensional-centred) object.
\end{proof}

\begin{theorem}[The commutant (self-intertwiner) von Neumann algebra]
    \label{thrm:commutant-von-neumann}
    \lean{Physicslib4.GNS.commutantVonNeumann, Physicslib4.GNS.mem_commutantVonNeumann_iff_intertwines, Physicslib4.GNS.isIrreducible_iff_commutantVonNeumann_eq_scalars}
    \leanfile{Physicslib4/GNS/Superselection.lean}
    \uses{def:irreducible-representation, lmm:bicommutant-of-selfadjoint-is-von-neumann}
    \leanok
    The commutant $\pi(A)'$ of a representation is packaged as a von Neumann algebra --- the algebra of self-intertwiners, i.e. the intertwiner/gauge algebra of $\pi$. Its underlying set is the centralizer of $\pi(A)$, which is a von Neumann algebra by \ref{lmm:bicommutant-of-selfadjoint-is-von-neumann} applied to the self-adjoint set $\pi(A)$; an operator lies in it exactly when it is a self-intertwiner of $\pi$, and $\pi$ is irreducible if and only if this algebra is trivial ($\pi(A)' = \mathbb{C}\cdot 1$) --- the von Neumann form of Schur's lemma.
\end{theorem}
\begin{proof}
    \leanok
    \uses{def:irreducible-representation, lmm:endomorphism-scalar, lmm:bicommutant-of-selfadjoint-is-von-neumann}
    The underlying set is a von Neumann algebra by \ref{lmm:bicommutant-of-selfadjoint-is-von-neumann}, since $\pi(A)$ is self-adjoint ($\pi$ being a $*$-homomorphism); membership unfolds to being a self-intertwiner of $\pi$, and the triviality criterion is \ref{lmm:endomorphism-scalar}.
\end{proof}

\begin{theorem}[Double-Commutant Duality]
    \label{thrm:double-commutant-duality}
    \lean{Physicslib4.GNS.commutant_gnsVonNeumannAlgebra, Physicslib4.GNS.commutant_commutantVonNeumann}
    \leanfile{Physicslib4/GNS/Superselection.lean}
    \uses{thrm:commutant-von-neumann}
    \leanok
    The generated von Neumann algebra $\pi(A)''$ and the commutant von Neumann algebra $\pi(A)'$ (\ref{thrm:commutant-von-neumann}) are each other's commutants.
\end{theorem}
\begin{proof}
    \leanok
    \uses{thrm:commutant-von-neumann}
    On one side, the commutant of $\pi(A)''$ is $\pi(A)'$: this is an instance of the triple-commutant collapse $S''' = S'$ applied to the self-adjoint image $\pi(A)$. On the other side, the commutant of $\pi(A)'$ is $\pi(A)''$, which is exactly the definition of the bicommutant. Together these are von Neumann's double-commutant relation for this pair, and they exhibit the bicommutant as idempotent: taking the commutant twice returns $(\pi(A)'')'' = \pi(A)''$, so $\pi(A)''$ and $\pi(A)'$ form a mutually dual pair under $S \mapsto S'$.
\end{proof}

\begin{theorem}[A Factor and its Commutant; Triviality Duality]
    \label{thrm:commutant-factor-duality}
    \lean{Physicslib4.GNS.isFactor_gnsVonNeumann_iff_isFactor_commutant, Physicslib4.GNS.commutantVonNeumann_eq_scalars_iff_gnsVonNeumann_eq_univ}
    \leanfile{Physicslib4/GNS/Superselection.lean}
    \uses{thrm:commutant-von-neumann}
    \leanok
    A von Neumann algebra and its commutant share the same center. Consequently $\pi(A)''$ is a factor (trivial center) if and only if its commutant $\pi(A)'$ is a factor. Dually, at the extreme of triviality, the commutant collapses to the scalars $\pi(A)' = \mathbb{C}\cdot 1$ if and only if the generated algebra is everything $\pi(A)'' = \mathcal{B}(H)$.
\end{theorem}
\begin{proof}
    \leanok
    \uses{thrm:double-commutant-duality, thrm:commutant-von-neumann, thrm:irreducible-generates-all}
    The intersection $\pi(A)'' \cap (\pi(A)'')' = \pi(A)'' \cap \pi(A)'$ is symmetric under the duality of \ref{thrm:double-commutant-duality}, being simultaneously the center of $\pi(A)''$ and the center of $\pi(A)'$. The triviality statement is the commutant form of the equivalence ``irreducible $\iff$ generates $\mathcal{B}(H)$'' (\ref{thrm:commutant-von-neumann}), passing between the two sides through the bicommutant $(\mathbb{C}\cdot 1)' = \mathcal{B}(H)$.
\end{proof}

\begin{lemma}[Abelian $\iff$ Self-Commuting]
    \label{lmm:von-neumann-abelian-self-commuting}
    \lean{Physicslib4.GNS.isAbelian_iff_le_commutant}
    \leanok
    A von Neumann algebra $R$ is abelian --- every pair of its elements commutes --- if and only if it is contained in its own commutant, $R \subseteq R'$. It is the operator-algebraic characterization of commutativity via the commutant, and the boundary case of microcausality $R(\mathbf{B}_1) \subseteq R(\mathbf{B}_2)'$ where the two regions coincide.
\end{lemma}
\begin{proof}
    \leanok
    This is essentially definitional: $R \subseteq R'$ says exactly that each element of $R$ commutes with every element of $R$.
\end{proof}

\begin{definition}[Center of a von Neumann algebra]
    \label{def:von-neumann-center}
    \lean{Physicslib4.GNS.vonNeumannCenter}
    \leanok
    The center of a von Neumann algebra $R$ on a Hilbert space $H$ is $Z(R) = R \cap R'$, the operators of $R$ that commute with every element of $R$. Since Mathlib's \texttt{VonNeumannAlgebra} carries no lattice meet, it is built as the meet of the star-subalgebras of $R$ and its commutant $R'$; its underlying set is $R \cap R'$. That this set is again a von Neumann algebra is recorded separately in \ref{lmm:von-neumann-inter-is-von-neumann}.
\end{definition}

\begin{lemma}[The center is a von Neumann algebra]
    \label{lmm:von-neumann-inter-is-von-neumann}
    \lean{Physicslib4.GNS.bicommutant_inter_commutant_eq}
    \uses{def:von-neumann-center}
    \leanok
    Let $R$ be a bundled von Neumann algebra on a Hilbert space $H$. Then $R \cap R'$ is bicommutant-closed, $(R \cap R')'' = R \cap R'$, so the center $Z(R) = R \cap R'$ of \ref{def:von-neumann-center} is again a von Neumann algebra.

    This is stated for the \emph{commutant pair} $M = R$, $N = R'$ only, because that is exactly what the Lean declaration proves: \texttt{bicommutant\_inter\_commutant\_eq} takes a single \texttt{VonNeumannAlgebra R} and forms $R \cap R^{\prime}$. The general two-algebra statement is \ref{lmm:von-neumann-inter-general}, which is not yet formalized.
\end{lemma}
\begin{proof}
    \leanok
    \uses{def:von-neumann-center, lmm:von-neumann-inter-general}
    The instance $M = R$, $N = R'$ of \ref{lmm:von-neumann-inter-general}.
\end{proof}

\begin{lemma}[The intersection of two von Neumann algebras is a von Neumann algebra]
    \label{lmm:von-neumann-inter-general}
    \lean{Physicslib4.GNS.bicommutant_inter_eq}
    \leanok
    Let $M$ and $N$ be bundled von Neumann algebras on a Hilbert space $H$. Then their underlying sets intersect in a bicommutant-closed set: $(M \cap N)'' = M \cap N$, so $M \cap N$ is again a von Neumann algebra. Here $M$ and $N$ are arbitrary and need \emph{not} form a commutant pair.

    This general form is what \ref{def:relative-commutant} needs, the relative commutant $R(\mathbf{B}_1)' \cap R(\mathbf{B}_2)$ being the instance $M = R(\mathbf{B}_1)'$, $N = R(\mathbf{B}_2)$, and likewise for its curved-spacetime counterpart \ref{def:relative-commutant-in-curved-spacetime}. It is formalized as \texttt{Physicslib4.GNS.bicommutant\_inter\_eq}, sitting alongside the commutant-pair instance \texttt{Physicslib4.GNS.bicommutant\_inter\_commutant\_eq} of \ref{lmm:von-neumann-inter-is-von-neumann}. One nuance is worth recording: in Lean the two are independent declarations, since \texttt{bicommutant\_inter\_commutant\_eq} proves the commutant-pair case directly rather than by specialising \texttt{bicommutant\_inter\_eq}, so when \ref{lmm:von-neumann-inter-is-von-neumann} describes itself as the instance $M = R$, $N = R'$ of this lemma that records the mathematics, not the Lean proof structure; the special case could be re-derived from the general one, but currently is not.
\end{lemma}
\begin{proof}
    \leanok
    Both $M$ and $N$ are bicommutant-closed, $M'' = M$ and $N'' = N$, so by the centralizer-of-a-union identity (Mathlib's \texttt{Set.centralizer\_union}) $M \cap N = M'' \cap N'' = (M' \cup N')'$. Thus $M \cap N$ is itself a commutant, and every commutant is bicommutant-closed by the triple-commutant collapse $S''' = S'$ (Mathlib's \texttt{Set.centralizer\_centralizer\_centralizer}): $(M \cap N)'' = (M' \cup N')''' = (M' \cup N')' = M \cap N$.
\end{proof}

\begin{theorem}[The center of a von Neumann algebra is abelian]
    \label{thrm:von-neumann-center-abelian}
    \lean{Physicslib4.GNS.vonNeumannCenter_isAbelian}
    \uses{def:von-neumann-center}
    \leanok
    The center $Z(R) = R \cap R'$ is abelian: any two of its elements commute.
\end{theorem}
\begin{proof}
    \leanok
    \uses{def:von-neumann-center, lmm:von-neumann-abelian-self-commuting}
    Indeed an element of the center lies in $R'$, hence commutes with every element of $R$, and in particular with every other element of the center (which lies in $R$).
\end{proof}

\begin{lemma}[$R$ is abelian iff it equals its center]
    \label{lmm:von-neumann-center-eq-self-iff-abelian}
    \lean{Physicslib4.GNS.vonNeumannCenter_eq_self_iff_isAbelian}
    \uses{def:von-neumann-center, lmm:von-neumann-abelian-self-commuting}
    \leanok
    A von Neumann algebra $R$ is abelian if and only if its center is all of $R$, i.e. $Z(R) = R \cap R' = R$. By \ref{lmm:von-neumann-abelian-self-commuting}, $R$ is abelian iff $R \subseteq R'$, and $R \cap R' = R$ holds exactly when $R \subseteq R'$.
\end{lemma}
\begin{proof}
    \leanok
    \uses{lmm:von-neumann-abelian-self-commuting}
    Abelianness of $R$ is equivalent to $R \subseteq R'$ (\ref{lmm:von-neumann-abelian-self-commuting}), which in turn is equivalent to $R \cap R' = R$.
\end{proof}

\begin{theorem}[A factor is abelian iff it is the scalars]
    \label{thrm:factor-abelian-iff-scalars}
    \lean{Physicslib4.GNS.isAbelian_iff_eq_scalars_of_isFactor}
    \uses{def:von-neumann-center}
    \leanok
    A factor $R$ (a von Neumann algebra with trivial center $R \cap R' = \mathbb{C}\cdot 1$) is abelian if and only if it equals the scalars $\mathbb{C}\cdot 1$.
\end{theorem}
\begin{proof}
    \leanok
    \uses{def:von-neumann-center, lmm:von-neumann-center-eq-self-iff-abelian, lmm:von-neumann-abelian-self-commuting}
    If $R$ is abelian then its center equals $R$ (\ref{lmm:von-neumann-center-eq-self-iff-abelian}), while by the factor hypothesis the center is $\mathbb{C}\cdot 1$; hence $R = \mathbb{C}\cdot 1$. Conversely the scalars are abelian, since $\mathbb{C}\cdot 1 \subseteq (\mathbb{C}\cdot 1)'$ (\ref{lmm:von-neumann-abelian-self-commuting}). Thus among factors, the abelian ones are exactly the trivial (one-dimensional) algebra of scalar multiples of the identity.
\end{proof}

\begin{theorem}[A von Neumann algebra and its commutant share a center]
  \label{thrm:center-eq-commutant-center}
  \lean{Physicslib4.GNS.vonNeumannCenter_eq_commutant}
  \uses{def:von-neumann-center}
  \leanok
  The center of a von Neumann algebra equals the center of its commutant: $Z(R) = Z(R')$.
\end{theorem}
\begin{proof}
  \leanok
  \uses{def:von-neumann-center, thrm:double-commutant-duality}
  Indeed $Z(R) = R \cap R'$ and $Z(R') = R' \cap R'' = R' \cap R$, which coincide by commutativity of intersection together with the double-commutant identity $R'' = R$.
\end{proof}

\begin{theorem}[A von Neumann algebra is a factor iff its center is the scalars]
  \label{thrm:factor-iff-center-scalars}
  \lean{Physicslib4.GNS.isFactor_iff_center_eq_scalars}
  \uses{def:von-neumann-center}
  \leanok
  A von Neumann algebra $R$ is a factor if and only if its center is the scalars, $Z(R) = R \cap R' = \mathbb{C}\cdot 1$.
\end{theorem}
\begin{proof}
  \leanok
  \uses{def:von-neumann-center}
  This is a restatement of the definition of a factor in terms of the bundled center: the underlying set of $Z(R)$ is exactly $R \cap R'$.
\end{proof}

\begin{theorem}[The Pure-State Dichotomy]
    \label{thrm:pure-state-dichotomy}
    \lean{Physicslib4.GNS.exists_gns_areDisjoint_or_unitaryEquiv_of_isPure}
    \leanfile{Physicslib4/GNS/Superselection.lean}
    \uses{def:pure-state}
    \leanok
    The GNS representations of two pure states are either disjoint or unitarily equivalent.
\end{theorem}
\begin{proof}
    \leanok
    \uses{thrm:irreducible-dichotomy, thrm:pure-iff-irreducible}
    Indeed the GNS representation of a pure state is irreducible (\ref{thrm:pure-iff-irreducible}), so the irreducible dichotomy (\ref{thrm:irreducible-dichotomy}) applies. Pure states thus fall into \emph{superselection sectors}: the sector of a pure state is the unitary-equivalence class of its irreducible GNS representation.
\end{proof}

\subsection{Direct Sums, Amplification, and Reducibility}

Given a family of $*$-representations $\pi_i : A \to \mathcal{B}(H_i)$, their direct sum acts on the $\ell^2$-direct sum $\bigoplus_i H_i$. This is the construction underlying amplifications and the reducibility of non-primary representations.

\begin{definition}[Direct-Sum Representation]
    \label{def:direct-sum-representation}
    \lean{Physicslib4.GNS.directSum}
    \leanfile{Physicslib4/GNS/DirectSum.lean}
    \leanok
    \uses{thrm:gns-construction-theorem}
    The \emph{direct sum} $\bigoplus_i \pi_i : A \to \mathcal{B}(\ell^2(\iota, H))$ of a family of $*$-representations acts coordinatewise: $(\bigoplus_i \pi_i)(a)$ is the diagonal operator $x \mapsto (\pi_i(a)\,x_i)_i$ on the $\ell^2$-direct sum. It is a $*$-representation, the diagonal being uniformly bounded by $\Vert a\Vert$ since each $\pi_i$ is contractive.
\end{definition}

\begin{theorem}[Subrepresentations and Commutant of a Direct Sum]
    \label{thrm:direct-sum-subrepresentation}
    \lean{Physicslib4.GNS.intertwines_single, Physicslib4.GNS.summandProj_mem_commutant}
    \leanfile{Physicslib4/GNS/DirectSum.lean}
    \uses{def:direct-sum-representation, def:disjoint-representations}
    \leanok
    Each summand embeds as a subrepresentation: the isometric inclusion $H_j \hookrightarrow \ell^2(\iota, H)$ intertwines $\pi_j$ with $\bigoplus_i \pi_i$. Moreover the orthogonal projection onto the $j$-th summand lies in the commutant of the direct sum, so each summand is a reducing subspace.
\end{theorem}
\begin{proof}
    \leanok
    \uses{def:direct-sum-representation, def:disjoint-representations}
    Both claims are coordinatewise computations with the diagonal action of \ref{def:direct-sum-representation}: the isometric inclusion is compatible with the $j$-th coordinate of the diagonal, and the projection onto the $j$-th summand commutes with every diagonal operator.
\end{proof}

\begin{definition}[Amplification]
    \label{def:amplification}
    \lean{Physicslib4.GNS.amplification}
    \leanfile{Physicslib4/GNS/Amplification.lean}
    \leanok
    \uses{def:direct-sum-representation}
    The \emph{$\iota$-fold amplification} $\iota \cdot \pi := \bigoplus_{i : \iota} \pi$ of a single representation is the direct sum of $\iota$ copies of $\pi$; $\pi$ embeds as each summand.
\end{definition}

\begin{theorem}[Reducibility of a Direct Sum]
    \label{thrm:direct-sum-reducible}
    \lean{Physicslib4.GNS.not_isIrreducible_directSum}
    \leanfile{Physicslib4/GNS/Amplification.lean}
    \uses{def:direct-sum-representation, def:irreducible-representation}
    \leanok
    A direct sum with two summands carrying nonzero vectors is \emph{reducible}. In particular an amplification $\iota \cdot \pi$ with at least two copies (and $\pi$ acting on a nonzero space) is reducible --- the multiplicity is visible in the commutant.
\end{theorem}
\begin{proof}
    \leanok
    \uses{def:irreducible-representation, thrm:direct-sum-subrepresentation}
    Were $\bigoplus_i \pi_i$ irreducible, every summand projection would be a scalar (its commutant being trivial); but the projection onto one of two nonzero summands is a nontrivial projection, hence not scalar.
\end{proof}

\subsection{Covariant States and the Covariance Action}

The Lorentz covariance of a Haag-Kastler net (\ref{def:lorentz-covariance}) acts on the net fiberwise, through the $*$-isomorphisms $\alpha_L : \mathfrak{U}(\mathbf{B}) \to \mathfrak{U}(L\cdot\mathbf{B})$. We record two consequences: the notion of a Lorentz-covariant family of local states, and the lift of the fiberwise action to a single dynamical $*$-automorphism of the quasilocal algebra.

\begin{definition}[Covariant Family of Local States]
    \label{def:covariant-state-family}
    \lean{Physicslib4.AQFT.HaagKastler.HaagKastlerNet.IsCovariantFamily}
    \leanfile{Physicslib4/AQFT/HaagKastler/CovariantState.lean}
    \leanok
    \uses{def:haag-kastler-net, def:state, def:lorentz-covariance}
    Given a Haag-Kastler net (\ref{def:haag-kastler-net}), a \emph{covariant family of local states} assigns to every region $\mathbf{B}$ a state $\omega_{\mathbf{B}}$ on the local algebra $\mathfrak{U}(\mathbf{B})$ such that, for every Lorentz transformation $L$ and every $a \in \mathfrak{U}(\mathbf{B})$, one has $\omega_{\mathbf{B}}(a) = \omega_{L\cdot\mathbf{B}}(\alpha_L a)$, where $\alpha_L$ is the covariance isomorphism of \ref{def:lorentz-covariance}. The local states are thus intertwined by the Lorentz action.
\end{definition}

\begin{lemma}[Composition of Covariance]
    \label{lmm:covariant-state-family-compose}
    \lean{Physicslib4.AQFT.HaagKastler.HaagKastlerNet.IsCovariantFamily.comp}
    \leanfile{Physicslib4/AQFT/HaagKastler/CovariantState.lean}
    \uses{def:covariant-state-family}
    \leanok
    For a covariant family of local states, the covariance relation composes along the group: $\omega_{\mathbf{B}}(a) = \omega_{L'\cdot(L\cdot\mathbf{B})}\big(\alpha_{L'}(\alpha_L a)\big)$.
\end{lemma}
\begin{proof}
    \leanok
    \uses{def:covariant-state-family}
    This reflects the multiplicativity of the Lorentz action: the covariance relation of \ref{def:covariant-state-family} is applied twice, first for $L$ and then for $L'$.
\end{proof}

\begin{definition}[Quasilocal Covariance Automorphism]
    \label{def:quasilocal-lift}
    \lean{Physicslib4.AQFT.HaagKastler.HaagKastlerNet.QuasilocalLift}
    \leanfile{Physicslib4/AQFT/HaagKastler/QuasilocalAction.lean}
    \leanok
    \uses{def:quasilocal-algebra, def:lorentz-covariance}
    A \emph{lift} of the fiberwise Lorentz action of $L$ to a quasilocal algebra $\mathfrak{U}$ (\ref{def:quasilocal-algebra}) is a $*$-automorphism $\beta_L$ of $\mathfrak{U}$ intertwining the local embeddings $\iota_{\mathbf{B}}$ with the covariance isomorphisms: $\beta_L(\iota_{\mathbf{B}} a) = \iota_{L\cdot\mathbf{B}}(\alpha_L a)$ for every Alexandrov-basis set $\mathbf{B}$.
\end{definition}

\begin{lemma}[Uniqueness of the Quasilocal Lift]
    \label{lmm:quasilocal-lift-unique}
    \lean{Physicslib4.AQFT.HaagKastler.HaagKastlerNet.QuasilocalLift.unique}
    \leanfile{Physicslib4/AQFT/HaagKastler/QuasilocalAction.lean}
    \uses{def:quasilocal-lift}
    \leanok
    Any two lifts of the same $L$ have the same underlying automorphism.
\end{lemma}
\begin{proof}
    \leanok
    \uses{def:quasilocal-lift}
    They agree on the union of the local images, which is dense in $\mathfrak{U}$, and $*$-automorphisms of a C*-algebra are continuous.
\end{proof}

\begin{theorem}[Existence of the Quasilocal Lift]
    \label{thrm:quasilocal-lift-exists}
    \lean{Physicslib4.AQFT.HaagKastler.nonempty_quasilocalLift}
    \leanfile{Physicslib4/AQFT/HaagKastler/QuasilocalIntertwiner.lean}
    \uses{def:quasilocal-lift, def:quasilocal-algebra}
    \leanok
    For a covariance-compatible quasilocal algebra, the fiberwise Lorentz action extends to a $*$-automorphism of $\mathfrak{U}$, so the lift of \ref{def:quasilocal-lift} exists.
\end{theorem}
\begin{proof}
    \leanok
    \uses{def:quasilocal-lift, def:quasilocal-algebra}
    The intertwiner is defined on the directed union of local images (a dense $*$-subalgebra) and extended by uniform continuity; the inverse is supplied by $L^{-1}$, and the group laws $\beta_{L'L} = \beta_{L'}\circ\beta_L$ and $\beta_{1} = \mathrm{id}$ furnish the two-sided inverse.
\end{proof}

\begin{theorem}[Existence for the Trivial Net]
    \label{thrm:quasilocal-lift-trivial}
    \lean{Physicslib4.AQFT.HaagKastler.nonempty_trivialQuasilocalLift}
    \leanfile{Physicslib4/AQFT/HaagKastler/QuasilocalIntertwiner.lean}
    \leanok
    The covariance-compatibility hypothesis is satisfiable: the trivial net's quasilocal algebra is covariance-compatible, so the quasilocal lift exists unconditionally for the trivial net.
\end{theorem}
\begin{proof}
    \leanok
    \uses{thrm:quasilocal-lift-exists}
    Covariance-compatibility holds because every $*$-automorphism of $\mathbb{C}$ is the identity; the lift is then supplied by \ref{thrm:quasilocal-lift-exists}.
\end{proof}

\begin{definition}[Covariant Quasilocal Algebra]
    \label{def:covariant-quasilocal-algebra}
    \lean{Physicslib4.AQFT.HaagKastler.CovariantQuasilocalAlgebra}
    \leanfile{Physicslib4/AQFT/HaagKastler/QuasilocalIntertwiner.lean}
    \leanok
    \uses{def:haag-kastler-net, def:quasilocal-algebra, def:lorentz-covariance, def:quasilocal-lift, thrm:quasilocal-lift-exists}
    A \emph{covariant quasilocal algebra} bundles a Haag-Kastler net (\ref{def:haag-kastler-net}), a quasilocal algebra of that net (\ref{def:quasilocal-algebra}), and a proof that the embeddings are covariance-compatible for every Lorentz transformation. On such data the quasilocal lift (\ref{def:quasilocal-lift}) exists for every $L$ by \ref{thrm:quasilocal-lift-exists}, yielding the covariance action $L \mapsto \beta_L$ on the quasilocal algebra; the trivial net provides an instance. This is the natural home for the covariance dynamics: the compatibility hypothesis of \ref{thrm:quasilocal-lift-exists} becomes structural data rather than a side condition.
\end{definition}

\begin{lemma}[Group-Action Coherence of the Covariance Automorphism]
    \label{lmm:quasilocal-action-coherence}
    \lean{Physicslib4.AQFT.HaagKastler.CovariantQuasilocalAlgebra.action_one, Physicslib4.AQFT.HaagKastler.CovariantQuasilocalAlgebra.action_mul}
    \leanfile{Physicslib4/AQFT/HaagKastler/QuasilocalIntertwiner.lean}
    \uses{def:covariant-quasilocal-algebra}
    \leanok
    The covariance action $L \mapsto \beta_L$ of a covariant quasilocal algebra is a genuine action of the Lorentz group by $*$-automorphisms of the quasilocal algebra: $\beta_{\mathbf{1}} = \mathrm{id}$ and $\beta_{L'L} = \beta_{L'} \circ \beta_L$.
\end{lemma}
\begin{proof}
    \leanok
    \uses{def:covariant-quasilocal-algebra}
    These are the group laws of the lifted covariance action carried by a covariant quasilocal algebra (\ref{def:covariant-quasilocal-algebra}).
\end{proof}

\begin{definition}[Invariant State]
    \label{def:invariant-state}
    \lean{Physicslib4.AQFT.HaagKastler.CovariantQuasilocalAlgebra.IsInvariantState}
    \leanfile{Physicslib4/AQFT/HaagKastler/QuasilocalIntertwiner.lean}
    \leanok
    \uses{def:covariant-quasilocal-algebra, def:state}
    A state $\omega$ on the quasilocal algebra of a covariant quasilocal algebra (\ref{def:covariant-quasilocal-algebra}) is \emph{(Poincar\'e-)invariant} if it is a fixed point of the dual covariance action: $\omega(\beta_L a) = \omega(a)$ for every Lorentz transformation $L$ and observable $a$. This is the invariance condition of a vacuum state; further conditions (e.g.\ the spectrum condition) are imposed separately.
\end{definition}

\begin{theorem}[GNS Unitary Implementation of an Invariant State]
    \label{thrm:invariant-state-gns-unitary}
    \lean{Physicslib4.AQFT.HaagKastler.CovariantQuasilocalAlgebra.IsInvariantState.exists_gns_unitary}
    \leanfile{Physicslib4/AQFT/HaagKastler/QuasilocalIntertwiner.lean}
    \uses{def:invariant-state, thrm:gns-construction-theorem}
    \leanok
    For an invariant state $\omega$, the covariance action is implemented on the GNS Hilbert space by a family of unitaries $U(L)$ satisfying $U(L)\,\pi(a)\Omega = \pi(\beta_L a)\,\Omega$ and $U(L)\Omega = \Omega$.
\end{theorem}
\begin{proof}
    \leanok
    \uses{def:invariant-state, lmm:quasilocal-action-coherence}
    The unitaries are obtained by extending the densely-defined isometry $\pi(a)\Omega \mapsto \pi(\beta_L a)\Omega$ - isometric because $\omega$ preserves the GNS inner product $\langle \pi(a)\Omega, \pi(b)\Omega\rangle = \omega(a^* b)$ - to the whole GNS space.
\end{proof}

Combining invariance with purity makes the GNS representation both covariant and irreducible. We emphasise that this is a \emph{precursor} to a vacuum representation, not a vacuum itself: a genuine vacuum additionally requires the spectrum condition (positivity of the energy-momentum spectrum), which is not imposed - and indeed not expressible - here.

\begin{theorem}[Irreducible Covariant Representation of a Pure Invariant State]
    \label{thrm:irreducible-covariant-representation}
    \lean{Physicslib4.AQFT.HaagKastler.CovariantQuasilocalAlgebra.IsInvariantState.exists_gns_irreducible_covariant}
    \leanfile{Physicslib4/AQFT/HaagKastler/QuasilocalIntertwiner.lean}
    \uses{def:pure-state, def:irreducible-representation, thrm:irreducible-generates-all}
    \leanok
    A state $\omega$ on the quasilocal algebra that is both invariant under the covariance action and pure yields a GNS representation that is simultaneously \emph{covariant} - implemented by a unitary representation $U(L)$ of the inhomogeneous Lorentz group fixing the cyclic vector $\Omega$, with operator covariance $U(L)\,\pi(a)\,U(L)^{-1} = \pi(\beta_L a)$ - and \emph{irreducible}; in particular it generates all of $\mathcal{B}(H)$ ($\pi(\mathfrak{U})'' = \mathcal{B}(H)$, by \ref{thrm:irreducible-generates-all}). It is a necessary precursor to a vacuum representation; the spectrum condition would be the remaining ingredient.
\end{theorem}
\begin{proof}
    \leanok
    \uses{thrm:invariant-state-gns-unitary, thrm:pure-iff-irreducible}
    This combines the covariant GNS triple of \ref{thrm:invariant-state-gns-unitary} with purity $\Rightarrow$ irreducibility (\ref{thrm:pure-iff-irreducible}).
\end{proof}

\begin{definition}[Positive Energy (bounded-generator scaffold)]
    \label{def:positive-energy}
    \lean{Physicslib4.AQFT.IsPositiveEnergy}
    \leanfile{Physicslib4/AQFT/HaagKastler/VacuumState.lean}
    \leanok
    A strongly continuous one-parameter unitary group $V : \mathbb{R} \to \mathcal{U}(H)$ has \emph{positive energy} when its generator is a positive operator: there exists a positive bounded operator $P$ (hence self-adjoint, with non-negative spectrum) such that $V(t) = e^{i t P}$ for all $t$. The positivity of $P$ is the energy-positivity asserted by the spectrum condition. The generator of a physical translation is unbounded, so requiring $P$ bounded is a genuine restriction; the faithful unbounded form requires Stone's theorem and the theory of unbounded self-adjoint operators, which Mathlib does not yet provide. This is the bounded-generator scaffold.
\end{definition}

\begin{theorem}[Positive-Energy API]
    \label{thrm:positive-energy-api}
    \lean{Physicslib4.AQFT.isPositiveEnergy_const_refl, Physicslib4.AQFT.exp_generator_unique, Physicslib4.AQFT.IsPositiveEnergy.conj, Physicslib4.AQFT.IsPositiveEnergy.strongContinuous}
    \leanfile{Physicslib4/AQFT/HaagKastler/VacuumState.lean}
    \uses{def:positive-energy}
    \leanok
    Structural properties of the positive-energy condition, all Stone-free. (i) \emph{Trivial subgroup}: the constant group $t \mapsto \mathrm{id}$ has positive energy, with zero generator. (ii) \emph{Uniqueness of the generator}: if two bounded generators induce the same one-parameter group, $e^{i t P} = e^{i t Q}$ for all $t$, then $P = Q$. Hence the generator witnessing positive energy is unique. (iii) \emph{Unitary invariance}: if $V$ has positive energy, so does its conjugate $t \mapsto W \circ V(t) \circ W^{-1}$ by a unitary $W$, with generator $W P W^{-1}$. Physically, the spectrum condition does not depend on the choice of unitary frame. (iv) \emph{Strong continuity}: a positive-energy group is strongly continuous --- $t \mapsto V(t) x$ is continuous for every $x$. This justifies calling it a strongly continuous one-parameter unitary group.
\end{theorem}
\begin{proof}
    \leanok
    \uses{def:positive-energy}
    (ii) Differentiating at $t = 0$ gives $i P = i Q$; positivity is not needed. (iii) The generator $W P W^{-1}$ is positive, being the unitary conjugate of $P$; the energy exponential transports via $e^{W A W^{-1}} = W e^{A} W^{-1}$. (iv) $t \mapsto (i t) P$ is continuous and the operator exponential is continuous.
\end{proof}

\begin{definition}[Vacuum State (generator-parameterized scaffold)]
    \label{def:vacuum-state}
    \lean{Physicslib4.AQFT.HaagKastler.CovariantQuasilocalAlgebra.IsVacuumState}
    \leanfile{Physicslib4/AQFT/HaagKastler/VacuumState.lean}
    \leanok
    \uses{def:invariant-state, def:positive-energy, def:covariant-quasilocal-algebra, thrm:invariant-state-gns-unitary}
    A state $\omega$ on the quasilocal algebra is a \emph{vacuum state} (relative to a future-timelike-translation predicate) when it is invariant under the covariance action (\ref{def:invariant-state}) and, in a GNS representation reproducing $\omega$ and implementing the action by unitaries $U(L)$, every future-timelike translation one-parameter subgroup $\gamma$ has positive energy: $t \mapsto U(\gamma(t))$ satisfies \ref{def:positive-energy}. This packages the two \emph{necessary} vacuum conditions --- invariance and the spectrum condition --- with the spectrum condition entering as the positive-energy hypothesis on the implementing unitaries. The future-timelike-translation predicate is a parameter, to be instantiated once the translation subgroup of the Lorentz group and its causal structure are wired in; the positive-energy condition is the bounded-generator scaffold. Constructing/discharging the spectrum condition for a concrete net is the Stone-gated next layer.
\end{definition}

\begin{theorem}[No-Stone Consequences of a Vacuum State]
    \label{thrm:vacuum-no-stone}
    \lean{Physicslib4.AQFT.HaagKastler.CovariantQuasilocalAlgebra.IsVacuumState.invariant, Physicslib4.AQFT.HaagKastler.CovariantQuasilocalAlgebra.IsVacuumState.exists_gns_irreducible_covariant}
    \leanfile{Physicslib4/AQFT/HaagKastler/VacuumState.lean}
    \uses{def:vacuum-state, thrm:irreducible-covariant-representation}
    \leanok
    The conditions a vacuum state satisfies that need neither the spectrum condition nor Stone's theorem. First, a vacuum state is \emph{invariant} (the first conjunct of the definition). Second, a \emph{pure} vacuum state yields the irreducible covariant GNS representation of \ref{thrm:irreducible-covariant-representation}: a covariant GNS triple with implementing unitaries $U(L)$ fixing $\Omega$ and operator covariance $U(L)\pi(a)U(L)^{-1} = \pi(\beta_L a)$, whose representation is irreducible and generates all of $\mathcal{B}(H)$.
\end{theorem}
\begin{proof}
    \leanok
    \uses{def:vacuum-state, thrm:irreducible-covariant-representation}
    Both follow by projecting to invariance and chaining with purity $\Rightarrow$ irreducibility; the spectrum-condition content of the vacuum is not used.
\end{proof}

\begin{definition}[Future-Timelike Translation Subgroup]
    \label{def:future-timelike-translation}
    \lean{Physicslib4.AQFT.HaagKastler.translationSub, Physicslib4.AQFT.HaagKastler.translationFlow, Physicslib4.AQFT.HaagKastler.isOneParameterSubgroup_translationFlow, Physicslib4.AQFT.HaagKastler.IsFutureTimelikeTranslation}
    \leanfile{Physicslib4/AQFT/HaagKastler/VacuumState.lean}
    \leanok
    \uses{def:lorentz-covariance, def:vacuum-state}
    The pure \emph{translations} of the inhomogeneous Lorentz group: an element is a pair $(\text{linear}, \text{translation})$, and a pure translation $\mathrm{translationSub}(n) = (\mathrm{id}, n)$ has trivial linear part, so $n \mapsto (\mathrm{id}, n)$ embeds the additive group of the spacetime carrier ($\mathrm{translationSub}(0) = 1$, $\mathrm{translationSub}(n+m) = \mathrm{translationSub}(n)\,\mathrm{translationSub}(m)$). The one-parameter translation flow in a direction $n$ is $\mathrm{translationFlow}(n)(t) = (\mathrm{id}, t \cdot n)$, which is a one-parameter subgroup. A one-parameter subgroup $\gamma$ is a \emph{future-timelike translation} when $\gamma = \mathrm{translationFlow}(n)$ for some future-pointing timelike $n$ --- i.e. $n$ lies in the forward Minkowski cone at the origin. This wires in the translation subgroup and its causal structure, giving the concrete predicate with which the abstract future-timelike-translation parameter of \ref{def:vacuum-state} is discharged.
\end{definition}

\begin{definition}[Vacuum State with the Concrete Spectrum Condition]
    \label{def:vacuum-state-concrete}
    \lean{Physicslib4.AQFT.HaagKastler.CovariantQuasilocalAlgebra.IsVacuumStateConcrete}
    \leanfile{Physicslib4/AQFT/HaagKastler/VacuumState.lean}
    \leanok
    \uses{def:vacuum-state, def:future-timelike-translation}
    The vacuum-state condition of \ref{def:vacuum-state} with its future-timelike-translation parameter fixed to the concrete predicate \ref{def:future-timelike-translation}: the spectrum condition is imposed on exactly the one-parameter translation subgroups $t \mapsto (\mathrm{id}, t \cdot n)$ with $n$ future-pointing timelike. The vacuum definition then depends on no free predicate. Invariance and (for a pure state) the irreducible covariant representation follow as in \ref{thrm:vacuum-no-stone}, since the concrete form unfolds to the parameterized one.
\end{definition}

\begin{theorem}[Purity is Covariance-Invariant]
    \label{thrm:purity-covariance-invariant}
    \lean{Physicslib4.GNS.isPure_precomp_iff, Physicslib4.AQFT.HaagKastler.CovariantQuasilocalAlgebra.isPure_precomp_action_iff}
    \leanfile{Physicslib4/GNS/ExtremeState.lean}
    \uses{def:pure-state, lmm:quasilocal-action-coherence}
    \leanok
    Purity is preserved by any $*$-automorphism: for $\Phi : \mathfrak{U} \xrightarrow{\sim} \mathfrak{U}$, the pullback state $\omega \circ \Phi$ is pure if and only if $\omega$ is. Applied to the covariance automorphism $\Phi = \beta_L$, this says purity of a state is a Lorentz-covariance-invariant property.
\end{theorem}
\begin{proof}
    \leanok
    \uses{def:pure-state}
    A dominated positive functional $\psi \le \omega \circ \Phi$ transports to $\psi \circ \Phi^{-1} \le \omega$, which purity sends to a scalar multiple of $\omega$; transporting back gives $\psi$ proportional to $\omega \circ \Phi$.
\end{proof}

\begin{theorem}[GNS covariance along the quasilocal action]
    \label{thrm:gns-covariance-quasilocal-action}
    \lean{Physicslib4.AQFT.HaagKastler.CovariantQuasilocalAlgebra.unitaryEquiv_gns_action, Physicslib4.AQFT.HaagKastler.CovariantQuasilocalAlgebra.isIrreducible_iff_gns_action, Physicslib4.AQFT.HaagKastler.CovariantQuasilocalAlgebra.isFactor_iff_gns_action}
    \uses{def:covariant-quasilocal-algebra, thrm:gns-covariance-unitary-equiv, thrm:gns-sector-transport, def:state-pullback}
    \leanok
    The lifted covariance automorphism $\beta_L$ of the quasilocal algebra $\mathfrak{U}$ is a $*$-automorphism, hence in particular a $*$-isomorphism, so the abstract GNS-covariance results apply to it. Let $\omega$ be a state on $\mathfrak{U}$. A cyclic representation of $\mathfrak{U}$ reproducing the pullback state $\omega \circ \beta_L$ is unitarily equivalent to $\pi_\omega \circ \beta_L$ (\ref{thrm:gns-covariance-unitary-equiv}), is irreducible exactly when $\pi_\omega$ is, and generates a factor exactly when $\pi_\omega$ does (\ref{thrm:gns-sector-transport}). Whereas \ref{thrm:gns-covariance-local} moves between the algebras of two different regions, here the algebra is fixed and the Lorentz group acts on it: the conclusion is that the superselection type of a \emph{global} state is a Lorentz invariant. Together with the invariance of purity (\ref{thrm:purity-covariance-invariant}) this says the entire sector classification of global states is Lorentz invariant.
\end{theorem}
\begin{proof}
    \leanok
    \uses{thrm:gns-covariance-unitary-equiv, thrm:gns-sector-transport, def:covariant-quasilocal-algebra}
    Instantiate \ref{thrm:gns-covariance-unitary-equiv} and \ref{thrm:gns-sector-transport} at $A := \mathfrak{U}$, $B := \mathfrak{U}$ and $\Phi := \beta_L$, the $*$-automorphism of the quasilocal algebra supplied by the covariance action (\ref{def:covariant-quasilocal-algebra}). Source and target algebra coincide, so no region-cast intervenes.
\end{proof}

\subsection{The Separating Vector of a Faithful State}

A faithful state has a sharper consequence at the level of the GNS construction than the faithfulness of its representation: its cyclic vector is also \emph{separating}. This is the basic datum of the modular (Tomita-Takesaki) theory of the associated von Neumann algebra.

\begin{theorem}[Separating Vector of a Faithful State]
    \label{thrm:separating-faithful}
    \lean{Physicslib4.GNS.separating_of_faithful, Physicslib4.GNS.exists_gns_separating}
    \leanfile{Physicslib4/GNS/Separating.lean}
    \uses{def:state, def:cyclic-vector, thrm:gns-construction-theorem}
    \leanok
    Let $\omega$ be a faithful state and $(\mathcal{H}_\omega, \pi_\omega, \Omega)$ its GNS triple. Then the cyclic vector $\Omega$ is \emph{separating} for $\pi_\omega(\mathfrak{U})$: if $\pi_\omega(a)\Omega = 0$ then $a = 0$. This holds in any representation reproducing a faithful state, not only the GNS one.
\end{theorem}
\begin{proof}
    \leanok
    \uses{def:cyclic-vector, thrm:gns-construction-theorem}
    Indeed $\pi_\omega(a)\Omega = 0$ gives $\omega(a^* a) = \langle \Omega, \pi_\omega(a^* a)\Omega\rangle = \langle \pi_\omega(a)\Omega, \pi_\omega(a)\Omega\rangle = 0$, and faithfulness forces $a = 0$.
\end{proof}

\subsection{The KMS Condition and Thermal Equilibrium}

Poincar\'e (or Killing-flow) invariance is only part of what singles out a physical equilibrium state. The Kubo-Martin-Schwinger (KMS) condition is the algebraic characterization of thermal equilibrium at inverse temperature $\beta$. It is phrased purely as an \emph{analyticity} statement about correlation functions, so - unlike the spectrum condition - it needs no unbounded-operator theory (no Stone theorem, no spectral measures). It is exactly the condition satisfied by the curved-spacetime examples of \ref{thrm:gns-unitary-stabilizer-in-curved-spacetime}: the Hartle-Hawking state on a black-hole exterior and the Gibbons-Hawking state in the de Sitter static patch are KMS for the relevant Killing flow.

\begin{definition}[One-Parameter Automorphism Group]
    \label{def:one-parameter-aut}
    \lean{Physicslib4.AQFT.IsOneParameterAut}
    \leanfile{Physicslib4/AQFT/KMS.lean}
    \leanok
    A family $\alpha : \mathbb{R} \to (A \simeq_{\star\mathrm{a}} A)$ of $*$-automorphisms of a C$^*$-algebra $A$ is a \emph{one-parameter group} if $\alpha_0 = \mathrm{id}$ and $\alpha_{s+t} = \alpha_s \circ \alpha_t$. This is the algebraic time evolution; in the curved-spacetime setting it is the automorphism group induced by a Killing flow.
\end{definition}

\begin{definition}[KMS State]
    \label{def:kms-state}
    \lean{Physicslib4.AQFT.IsKMSState}
    \leanfile{Physicslib4/AQFT/KMS.lean}
    \leanok
    \uses{def:one-parameter-aut, def:state}
    A state $\omega$ on $A$ is \emph{$(\alpha, \beta)$-KMS} for a one-parameter automorphism group $\alpha$ (\ref{def:one-parameter-aut}) at inverse temperature $\beta$ if for every $a, b \in A$ the correlation function $t \mapsto \omega(a\,\alpha_t b)$ extends to a function $F$ on the closed strip $0 \le \operatorname{Im} z \le \beta$ that is continuous there, holomorphic on the open strip, bounded, and whose boundary value on $\operatorname{Im} z = \beta$ is $t \mapsto \omega(\alpha_t b\, \cdot a)$.
\end{definition}

\begin{theorem}[The KMS State Set is Convex]
    \label{thrm:kms-convex}
    \lean{Physicslib4.AQFT.IsKMSState.convexCombo}
    \leanfile{Physicslib4/AQFT/KMS.lean}
    \uses{def:kms-state}
    \leanok
    A convex combination $s\,\omega_1 + (1-s)\,\omega_2$ ($0 \le s \le 1$) of two $(\alpha, \beta)$-KMS states is again $(\alpha, \beta)$-KMS. Thus the equilibrium states at a fixed temperature form a convex set.
\end{theorem}
\begin{proof}
    \leanok
    \uses{def:kms-state}
    For each pair $(a, b)$ the analytic interpolant is the convex combination $s\,F_1 + (1-s)\,F_2$ of the two interpolants, which is continuous, bounded, and holomorphic on the strip, with boundary values that add.
\end{proof}

\begin{lemma}[Boundary Coincidence for $a = 1$]
    \label{lmm:kms-correlation-one}
    \lean{Physicslib4.AQFT.IsKMSState.correlationOne}
    \leanfile{Physicslib4/AQFT/KMS.lean}
    \uses{def:kms-state}
    \leanok
    For a KMS state $\omega$ and any observable $a$, the correlation function $F$ of the pair $(1, a)$ has its two boundary values \emph{equal} - both are $t \mapsto \omega(\alpha_t a)$.
\end{lemma}
\begin{proof}
    \leanok
    \uses{def:kms-state}
    This is the algebraic heart of the invariance argument: it follows directly from the KMS condition with $a := 1$, since $\omega(1\cdot\alpha_t a) = \omega(\alpha_t a \cdot 1) = \omega(\alpha_t a)$.
\end{proof}

\begin{definition}[Strip-Liouville Principle]
    \label{def:strip-liouville}
    \lean{Physicslib4.AQFT.StripLiouville}
    \leanfile{Physicslib4/AQFT/KMS.lean}
    \leanok
    The \emph{strip-Liouville principle} at width $\beta$ is the statement that any function $F$ continuous and bounded on the closed strip $0 \le \operatorname{Im} z \le \beta$, holomorphic on the open strip, and with equal boundary values $F(t) = F(t + i\beta)$ for all real $t$, is constant along the real axis: $F(t) = F(0)$. It is the analytic input that turns boundary coincidence into invariance.
\end{definition}

\begin{theorem}[$i\beta$-Periodic Entire Extension (Strip Schwarz Reflection)]
    \label{thrm:strip-periodic-extension}
    \lean{Physicslib4.exists_bounded_entire_extension_of_strip_periodic}
    \leanfile{Physicslib4/Analysis/StripPeriodicExtension.lean}
    \leanok
    A function $F$ continuous on the closed strip $0 \le \operatorname{Im} z \le \beta$, holomorphic on the open strip, bounded, and with equal boundary values $F(t) = F(t + i\beta)$ on the real axis, admits a \emph{bounded entire} extension $H$ agreeing with $F$ on $\mathbb{R}$. This is the analytic engine behind the strip-Liouville principle.
\end{theorem}
\begin{proof}
    \leanok
    The extension is the $i\beta$-periodic continuation $H(z) = F\!\left(z - \lfloor \operatorname{Im} z / \beta \rfloor\, i\beta\right)$, which folds every point into the fundamental strip: it is continuous across each gluing line $\operatorname{Im} z = k\beta$ because the boundary values match by periodicity, holomorphic off the lines as $F$ composed with a holomorphic shift, and holomorphic on the lines by the horizontal-line removable-singularity theorem (Morera).
\end{proof}

\begin{theorem}[Strip-Liouville Holds for $\beta > 0$]
    \label{thrm:strip-liouville-pos}
    \lean{Physicslib4.AQFT.stripLiouville_of_pos}
    \leanfile{Physicslib4/AQFT/KMS.lean}
    \uses{def:strip-liouville}
    \leanok
    At positive width $\beta > 0$ the strip-Liouville principle (\ref{def:strip-liouville}) is a theorem. The hypothesis $\beta > 0$ is necessary: at $\beta = 0$ the open strip is empty and at $\beta < 0$ the strip itself is empty, and in both cases the principle is false.
\end{theorem}
\begin{proof}
    \leanok
    \uses{def:strip-liouville, thrm:strip-periodic-extension}
    The equal boundary values let $F$ extend to a bounded $i\beta$-periodic \emph{entire} function by the strip Schwarz reflection (\ref{thrm:strip-periodic-extension}), which is then constant on $\mathbb{R}$ by Liouville's theorem.
\end{proof}

\begin{theorem}[KMS States are Invariant]
    \label{thrm:kms-invariance}
    \lean{Physicslib4.AQFT.IsKMSState.invariant_of_pos}
    \leanfile{Physicslib4/AQFT/KMS.lean}
    \uses{def:kms-state}
    \leanok
    At positive inverse temperature $\beta > 0$, every $(\alpha, \beta)$-KMS state $\omega$ is $\alpha$-invariant: $\omega(\alpha_t a) = \omega(a)$ for all $t$ and $a$.
\end{theorem}
\begin{proof}
    \leanok
    \uses{def:kms-state, lmm:kms-correlation-one, thrm:strip-liouville-pos}
    By boundary coincidence (\ref{lmm:kms-correlation-one}) the KMS correlation function of $(1, a)$ has equal boundary values $F(t) = F(t + i\beta) = \omega(\alpha_t a)$; the strip-Liouville principle for $\beta > 0$ (\ref{thrm:strip-liouville-pos}) forces $F(t) = F(0)$, so $\omega(\alpha_t a) = F(t) = F(0) = \omega(\alpha_0 a) = \omega(a)$.
\end{proof}

\begin{theorem}[Uniqueness on the Strip from Boundary Values]
    \label{thrm:strip-uniqueness}
    \lean{Physicslib4.eqOn_strip_of_eq_boundary}
    \leanfile{Physicslib4/Analysis/StripPeriodicExtension.lean}
    \leanok
    Two functions continuous and bounded on the closed strip $0 \le \operatorname{Im} z \le \beta$ (with $\beta > 0$), holomorphic on the open strip, that agree on \emph{both} boundary lines $\operatorname{Im} z = 0$ and $\operatorname{Im} z = \beta$, agree on the whole strip.
\end{theorem}
\begin{proof}
    \leanok
    \uses{thrm:strip-periodic-extension}
    Indeed their difference has vanishing boundary values, so its $i\beta$-periodic extension (\ref{thrm:strip-periodic-extension}) is a bounded entire function, constant by Liouville, equal to its value $0$ at the origin.
\end{proof}

\begin{theorem}[Uniqueness of the KMS Correlation Function]
    \label{thrm:kms-correlation-unique}
    \lean{Physicslib4.AQFT.IsKMSState.correlation_eqOn}
    \leanfile{Physicslib4/AQFT/KMS.lean}
    \uses{def:kms-state}
    \leanok
    For $\beta > 0$, the analytic completion of a KMS correlation function is unique: any two functions satisfying the KMS analytic data for the same pair $(a, b)$ - continuous and bounded on the strip, holomorphic on the interior, with the prescribed boundary values $t \mapsto \omega(a\,\alpha_t b)$ and $t \mapsto \omega(\alpha_t b\,\cdot a)$ - agree on the whole strip.
\end{theorem}
\begin{proof}
    \leanok
    \uses{def:kms-state, thrm:strip-uniqueness}
    This is the strip uniqueness (\ref{thrm:strip-uniqueness}) applied to the two analytic completions, which share both boundary values.
\end{proof}

\subsection{KMS States for the Covariance Flow}

A one-parameter subgroup $t \mapsto L_t$ of the inhomogeneous Lorentz group - for instance the time-translation subgroup - induces, via the quasilocal lift $\beta_L$ (\ref{lmm:quasilocal-action-coherence}), a one-parameter group of $*$-automorphisms of the global quasilocal algebra $\mathfrak{U}$. Unlike the curved case, where the absence of a global algebra forces a restriction to the stabilizer $\mathrm{Stab}(\mathbf{B})$, here $\beta_L$ is a genuine automorphism of the single algebra $\mathfrak{U}$ for every $L$, so no restriction is needed.

\begin{definition}[Covariance-Flow Automorphism Family]
    \label{def:flow-aut-covariance}
    \lean{Physicslib4.AQFT.HaagKastler.CovariantQuasilocalAlgebra.flowAut}
    \leanfile{Physicslib4/AQFT/HaagKastler/QuasilocalKMS.lean}
    \leanok
    \uses{def:covariant-quasilocal-algebra, lmm:quasilocal-action-coherence}
    Given a one-parameter subgroup $t \mapsto L_t$ of the inhomogeneous Lorentz group, the \emph{covariance-flow automorphism family} of $\mathfrak{U}$ is $t \mapsto \beta_{L_t}$, the covariance action evaluated along the flow.
\end{definition}

\begin{lemma}[A Lorentz One-Parameter Subgroup Induces a One-Parameter Group]
    \label{lmm:one-parameter-aut-flow-covariance}
    \lean{Physicslib4.AQFT.HaagKastler.CovariantQuasilocalAlgebra.isOneParameterAut_flowAut}
    \leanfile{Physicslib4/AQFT/HaagKastler/QuasilocalKMS.lean}
    \uses{def:flow-aut-covariance, def:one-parameter-aut}
    \leanok
    If $t \mapsto L_t$ is a one-parameter subgroup ($L_0 = 1$, $L_{s+t} = L_s L_t$), then the induced automorphisms $t \mapsto \beta_{L_t}$ of $\mathfrak{U}$ form a one-parameter automorphism group (\ref{def:one-parameter-aut}).
\end{lemma}
\begin{proof}
    \leanok
    \uses{def:flow-aut-covariance, lmm:quasilocal-action-coherence}
    This holds by the coherence of the covariance action (\ref{lmm:quasilocal-action-coherence}).
\end{proof}

\begin{definition}[KMS State for the Covariance Flow]
    \label{def:kms-state-for-flow-covariance}
    \lean{Physicslib4.AQFT.HaagKastler.CovariantQuasilocalAlgebra.IsKMSStateForFlow}
    \leanfile{Physicslib4/AQFT/HaagKastler/QuasilocalKMS.lean}
    \leanok
    \uses{lmm:one-parameter-aut-flow-covariance, def:kms-state, def:state}
    A state $\omega$ on the quasilocal algebra $\mathfrak{U}$ is a \emph{KMS state for the covariance flow} $L$ at inverse temperature $\beta$ if it satisfies the KMS condition (\ref{def:kms-state}) for the induced one-parameter automorphism group (\ref{lmm:one-parameter-aut-flow-covariance}).
\end{definition}

\begin{theorem}[Convexity of the Covariance-Flow KMS States]
    \label{thrm:kms-convex-for-flow-covariance}
    \lean{Physicslib4.AQFT.HaagKastler.CovariantQuasilocalAlgebra.IsKMSStateForFlow.convexCombo}
    \leanfile{Physicslib4/AQFT/HaagKastler/QuasilocalKMS.lean}
    \uses{def:kms-state-for-flow-covariance}
    \leanok
    A convex combination $s\,\omega_1 + (1-s)\,\omega_2$ ($0 \le s \le 1$) of two KMS states on $\mathfrak{U}$ for the same covariance flow $L$ at the same inverse temperature $\beta$ is again a KMS state for that flow. Physically, the equilibrium states for a one-parameter symmetry flow form a convex set.
\end{theorem}
\begin{proof}
    \leanok
    \uses{def:kms-state-for-flow-covariance, thrm:kms-convex}
    This specializes the abstract KMS convexity (\ref{thrm:kms-convex}) to the induced one-parameter group $\mathrm{flowAut}$.
\end{proof}

\begin{definition}[Ground State for a Covariance Flow]
    \label{def:ground-state-for-flow-covariance}
    \lean{Physicslib4.AQFT.HaagKastler.CovariantQuasilocalAlgebra.IsGroundStateForFlow, Physicslib4.AQFT.HaagKastler.CovariantQuasilocalAlgebra.IsGroundStateForFlow.invariant, Physicslib4.AQFT.HaagKastler.CovariantQuasilocalAlgebra.IsGroundStateForFlow.exists_strongContinuous_unitary}
    \leanfile{Physicslib4/AQFT/HaagKastler/QuasilocalKMS.lean}
    \leanok
    \uses{def:kms-state-for-flow-covariance, def:positive-energy, def:covariant-quasilocal-algebra}
    A state $\omega$ on the quasilocal algebra $\mathfrak{U}$ is a \emph{ground state} for a one-parameter subgroup $t \mapsto L_t$ of the inhomogeneous Lorentz group (e.g.\ a translation or boost flow) when it is invariant under the flow and, in a GNS representation reproducing $\omega$ and implementing the flow by unitaries $U(t)$ fixing $\Omega$, the one-parameter unitary group $t \mapsto U(t)$ has positive energy (\ref{def:positive-energy}). This is the ground-state ($\beta \to \infty$, spectrum-condition) counterpart of the covariance-flow KMS state \ref{def:kms-state-for-flow-covariance}: the stationary state whose flow generator, the Hamiltonian for a timelike flow, is positive. Two Stone-free consequences: a ground state is flow-invariant, and its implementing unitary group is strongly continuous. It is the Minkowski analogue of the curved Killing-flow ground state \ref{def:ground-state-for-flow-in-curved-spacetime}.
\end{definition}
